% Source A: FIFTH_REFERENCE.md; SHA-256 67e3b0220ce7173e0223b7f00fbf67ed5adef5fd0c83bf85c79015f84cb660d2
% Source B: HEAT_AND_COMPLEMENT.md; SHA-256 250e567605d2d1e508aad21c4942a7c028a5d82b468c089d4e2c9a9b6e78dd2c
\documentclass[11pt,a4paper]{article}
\usepackage[margin=25mm,headheight=14pt]{geometry}
\usepackage{fontspec}
\setmainfont{Latin Modern Roman}
\setsansfont{Latin Modern Sans}
\setmonofont{Latin Modern Mono}[Scale=0.86]
\usepackage{amsmath,amssymb}
\usepackage{unicode-math}
\setmathfont{Latin Modern Math}
\usepackage{microtype,graphicx,adjustbox}
\usepackage{longtable,booktabs,array}
\newcounter{none}
\usepackage{calc,xcolor}
\definecolor{linkblue}{HTML}{244B70}
\usepackage{fvextra}
\DefineVerbatimEnvironment{verbatim}{Verbatim}{fontsize=\small,breaklines=true,breakanywhere=true,breaksymbolleft={},breaksymbolright={},xleftmargin=5pt}
\usepackage{enumitem}
\setlist{itemsep=2pt,parsep=1pt,topsep=4pt}
\usepackage{fancyhdr}
\pagestyle{fancy}
\fancyhf{}
\fancyhead[L]{\small SU(2) Yang-Mills: the fifth reference}
\fancyhead[R]{\small 21 September 2026}
\fancyfoot[C]{\thepage}
\renewcommand{\headrulewidth}{0.3pt}
\usepackage{hyperref}
\hypersetup{colorlinks=true,linkcolor=linkblue,urlcolor=linkblue,pdftitle={The
fifth reference in SU(2) lattice Yang-Mills
theory},pdfauthor={KokunoYumeto},pdfsubject={Fifth-reference source correction, heat estimates and complete physical complement},pdfcreator={Pandoc and XeLaTeX}}
\usepackage{bookmark}
\urlstyle{same}
\setlength{\emergencystretch}{4em}
\hfuzz=0.2pt
\setlength{\parindent}{0pt}
\setlength{\parskip}{5pt plus 1pt minus 1pt}
\setcounter{secnumdepth}{2}
\setcounter{tocdepth}{2}
\providecommand{\tightlist}{\setlength{\itemsep}{0pt}\setlength{\parskip}{0pt}}
\providecommand{\pandocbounded}[1]{#1}
\let\originalsection\section
\renewcommand{\section}{\clearpage\originalsection}
\begin{document}
\begin{titlepage}
\vspace*{14mm}
{\sffamily\Large MATHEMATICAL READER\par}
\vspace{12mm}
{\LARGE\bfseries The fifth reference in SU(2) lattice Yang-Mills
theory\par}
\vspace{9mm}
{\large A convergent source correction, enlarged heat bounds, and the
full physical complement\par}
\vspace{13mm}
{\large KokunoYumeto\par}
\vspace{3mm}
{21 September 2026\par}
\vfill
\begin{minipage}{0.96\textwidth}
A mathematical introduction followed by both complete proof manuscripts. Original link coordinates, Haar pairings, scalar terms, physical time and the full complementary space are retained.

\medskip
The estimates hold on specified strong-coupling domains and in the spatial limit at fixed lattice spacing. A four-dimensional continuum mass gap is not claimed.

\medskip
The distinct earlier Jacobian and \(S^6\) branches credit Levent Alpöge. Their sources and the direct analytic antecedents of this continuation are identified separately.
\end{minipage}
\vspace{10mm}
\end{titlepage}
\pagenumbering{roman}
\tableofcontents
\clearpage
\pagenumbering{arabic}
\section{Reading the fifth-reference
continuation}\label{reading-the-fifth-reference-continuation}

The fifth reference is an explicitly calculated approximation to the
logarithm of the interacting vacuum. Its purpose is to control the whole
nonlinear remainder and then recover statements about the original
Hamiltonian: a physical spectral gap, heat correlations, and the energy
change when an observed family of states can mix with every other
physical state. The two manuscripts collected here carry out that
programme on a stated strong-coupling domain, uniformly in spatial
volume.

The source construction reaches the endpoint
\(\alpha_5\approx0.018424953576117616681\), corresponding to
\(g^2\geq1/(2\sqrt{\alpha_5})\approx3.683551984\). Its heat estimate
uses the strictly interior circle \(R=1/55\). At the stronger benchmark
\(g^2\geq13\), relaxation into the complete physical complement changes
the original energy and state metrics by less than one part in two
thousand; at \(g^2\geq16\), both changes are less than one part in
twenty-five thousand. These are estimates at fixed lattice spacing. No
four-dimensional continuum Yang-Mills construction or continuum mass gap
follows from them.

This introduction explains the objects, domains and uses of the results.
Appendix A reproduces the complete fifth-reference proof, with its
original labels F1-F42. Appendix B reproduces the complete heat and
complement proof, with labels H1-H34. The introduction writes
\(\alpha_5\) for the endpoint called \(\alpha\) in those sources. Their
full mathematical text and citations are retained, including statements
of what is not established.

\subsection{The original Hamiltonian and its physical
pairing}\label{the-original-hamiltonian-and-its-physical-pairing}

Fix an integer \(L\geq2\). The vertices are \(\{-L,\ldots,L\}^3\), the
links are all contained positive nearest-neighbour edges, and
\(\mathsf P_L\) is the set of contained elementary plaquettes. A link
carries \(U_e\in\mathrm{SU}(2)\). Physical functions are invariant under
the gauge transformations at every original vertex, including boundary
vertices. Their Hilbert space is the invariant subspace of
product-Haar-probability \(L^2\).

The elementary observable and original derivatives are

\begingroup\small\begin{equation*}
\adjustbox{max width=\linewidth}{$\displaystyle W_p=\operatorname{tr}\!\left(U_i(n)U_j(n+e_i)U_i(n+e_j)^{-1}U_j(n)^{-1}\right),
\qquad T_a=-i\sigma_a/2,$}
\end{equation*}\endgroup

\begingroup\small\begin{equation*}
\adjustbox{max width=\linewidth}{$\displaystyle X_{e,a}f=\left.\frac{\partial}{\partial t}f(\ldots,e^{tT_a}U_e,\ldots)\right|_{t=0},
\qquad K=-\sum_{e,a}X_{e,a}^2.$}
\end{equation*}\endgroup

Writing \(M=|\mathsf P_L|\) and \(S=\sum_pW_p\), the Hamiltonian is

\begingroup\small\begin{equation*}
\adjustbox{max width=\linewidth}{$\displaystyle H_L=\kappa K+\kappa\xi(2M-S),\qquad
\kappa=\frac{2g^2}{a},\qquad \xi=\frac1{4g^4},\qquad a,g>0.$}
\end{equation*}\endgroup

Its operator and form domains are the physical parts of \(H^2\) and
\(H^1\) on the compact product of link groups. Let \(\psi\) be its
positive unit ground vector, \(E_0\) its energy, \(\rho=\psi^2\), and
\(A=H_L-E_0\). The physical energy identity is

\begingroup\small\begin{equation*}
\adjustbox{max width=\linewidth}{$\displaystyle q_A(\psi f)=\kappa\int\rho\sum_{e,a}|X_{e,a}f|^2\,dU.$}
\end{equation*}\endgroup

This identity retains the interacting weight \(\rho\) and the original
Haar measure. The coefficient norms used to construct \(\psi\) are
auxiliary analytic tools; they do not replace this pairing.

The exponential vacuum and character/Casimir framework has a human
antecedent in Schütte, Zheng Weihong and Hamer {[}1{]}, especially
Sections 2-5. Appendix A states the conversion from their convention:
\(g_C=2^{3/4}g\), \(H_{\mathrm{ours}}=\sqrt2H_C+2\kappa\xi MI\), and
their commutator derivative is \(iX\).

\subsection{Why the fifth coefficient improves the source
domain}\label{why-the-fifth-coefficient-improves-the-source-domain}

Let \(P_Hf=\int f\,dU\), \(Q_H=I-P_H\), and define

\begingroup\small\begin{equation*}
\adjustbox{max width=\linewidth}{$\displaystyle \Gamma(f,h)=\sum_{e,a}(X_{e,a}f)(X_{e,a}h),\qquad
B(f,h)=K^{-1}Q_H\Gamma(f,h).$}
\end{equation*}\endgroup

The logarithmic source satisfies \(v=\xi v_1+B(v,v)\), with \(v_1=S/3\).
Each Fourier coefficient carries its original finite edge support. The
local coefficient norm weights a physical representation block by its
Casimir \(c_j=\sum_ej_e(j_e+1)\); the two additional bounds \(m(v)\) and
\(t(v)\) weight its total spin and its spin at a specified edge. Their
precise definitions are F4-F6. Gauge singlet constraints give
\(c_j\geq6j_e\) and \(c_j\geq3\) for a nonconstant physical block.
Consequently

\begingroup\small\begin{equation*}
\adjustbox{max width=\linewidth}{$\displaystyle \|B(f,h)\|_{\mathrm{loc}}
\leq3\{m(f)t(h)+m(h)t(f)\}
\leq\frac23\|f\|_{\mathrm{loc}}\|h\|_{\mathrm{loc}}.$}
\end{equation*}\endgroup

The compact-group coefficient argument proves the multiplication and
derivative estimates directly; Eymard {[}2{]} is its Fourier-algebra
antecedent. The finite calculation then improves the inputs to this
estimate by resolving original signed spin channels before taking
absolute values.

The preserved fifth catalogue has 662 representatives, 124,864 anchored
multisets, and 14,063 rational trace-monomial terms. On 321
representatives, every face is distinct and no link occurs more than
twice. There the signed rational recurrence F16 resolves all admissible
spin-zero and spin-one channels on repeated links. Complete polynomial
identities compare the channel formula with the original fifth
coefficient. On the other 341 representatives, F22-F24 use the
applicable lower-order channel bounds and full trace expressions. In
particular, the distinct-face recurrence is not applied to links with
three or more occurrences.

All 662 rows return to the original marked edge through their saved
signed coordinate maps. The resulting uniform upward bounds are

\begingroup\small\begin{equation*}
\adjustbox{max width=\linewidth}{$\displaystyle \|v_5\|_{\mathrm{loc}}<2476866,\qquad
m(v_5)<1638684,\qquad t(v_5)<190128.$}
\end{equation*}\endgroup

F25 gives the stronger exact rational bounds. F26 lists the five pairs
\((m_i,t_i)\) used in the endpoint calculation. The finite dual
certificates bound actual coefficient objectives; they do not assert
that the physical coefficient masses attain a polytope maximum.

\subsection{A complete correction, including the
endpoint}\label{a-complete-correction-including-the-endpoint}

Set \(q_5=\sum_{i=1}^5\xi^iv_i\) and \(R_5=\xi v_1+B(q_5,q_5)-q_5\). The
residual includes every term of degrees six through ten; F27 displays it
in full. With \(x=|\xi|\), define

\begingroup\small\begin{equation*}
\adjustbox{max width=\linewidth}{$\displaystyle b_{ij}=3(m_it_j+m_jt_i),\qquad
\ell(x)=\sum_{i=1}^5(m_i+4t_i)x^i,$}
\end{equation*}\endgroup

\begingroup\small\begin{equation*}
\adjustbox{max width=\linewidth}{$\displaystyle \delta(x)=\sum_{\substack{1\leq i,j\leq5\\i+j\geq6}}b_{ij}x^{i+j},
\qquad \mathscr D(x)=(1-\ell(x))^2-\frac83\delta(x).$}
\end{equation*}\endgroup

The first positive root of \(\mathscr D\) is \(\alpha_5\). The exact
enclosures are

\begingroup\small\begin{equation*}
\adjustbox{max width=\linewidth}{$\displaystyle 0.018424953576117616681<\alpha_5<0.018424953576117616682,$}
\end{equation*}\endgroup

\begingroup\small\begin{equation*}
\adjustbox{max width=\linewidth}{$\displaystyle 3.683551983985727304439<\frac1{2\sqrt{\alpha_5}}
<3.683551983985727304440.$}
\end{equation*}\endgroup

For \(|\xi|\leq\alpha_5\), the actual source linearization
\(L_5=I-2B(q_5,\cdot)\) has a convergent Neumann inverse. The correction
\(w=v-q_5\) solves \(L_5w=R_5+B(w,w)\) and obeys

\begingroup\small\begin{equation*}
\adjustbox{max width=\linewidth}{$\displaystyle \|w\|_{\mathrm{loc}}\leq w_*(x)
=\frac34\left(1-\ell(x)-\sqrt{\mathscr D(x)}\right).$}
\end{equation*}\endgroup

The Catalan expansion converges at \(x=\alpha_5\) as well. Its explicit
omitted-order estimate F32 decays at least as \((N+1)^{-1/2}\) at the
endpoint. This is why the closed source disk includes its boundary even
though a strict contraction estimate would not suffice there.

For real coupling, \(\psi=e^{v+c_L}\) with
\(c_L=-\frac12\log\int e^{2v}dU\) gives the original positive ground
vector. The same correction yields

\begingroup\small\begin{equation*}
\adjustbox{max width=\linewidth}{$\displaystyle \Delta_L\geq\kappa d_5(\xi),\qquad
d_5(x)=\frac32\left(1+\sqrt{\mathscr D(x)}\right)
+\sum_{i=1}^5\left(\frac32m_i-6t_i\right)x^i.$}
\end{equation*}\endgroup

Here \(\Delta_L\) is the gap of the whole physical Hamiltonian above its
vacuum. F35-F37 show \(d_5(x)>0\) throughout the closed interval and
prove the corresponding weighted Poincaré inequality on the full
physical form domain. For example, \(g^2\geq4\) gives
\(\Delta_L/\kappa>1.9068\) and \(\|v-q_5\|_{\mathrm{loc}}<0.005752\).
The comparison between the approximate reference and the physical
operator retains its scalar, multiplication and drift defects in
F39-F40.

\subsection{The enlarged heat circle and what its remainder
controls}\label{the-enlarged-heat-circle-and-what-its-remainder-controls}

For the centered plaquette states
\(r_p=(W_p-\langle W_p\rangle_\rho)\psi\), write

\begingroup\small\begin{equation*}
\adjustbox{max width=\linewidth}{$\displaystyle \widehat C_{pq}(\tau;\xi)=\langle r_p,e^{-\tau A/\kappa}r_q\rangle,
\qquad \tau=\kappa t.$}
\end{equation*}\endgroup

The physical time is \(t\). Appendix B constructs the full coefficient
generator, its dense domain and its range, and the complex mean with its
nonzero denominator. These are the hypotheses used for the
Lumer-Phillips generation theorem {[}3{]}; a bound on a single inverse
is not substituted for a heat estimate.

On the circle \(R=1/55\), the fifth-source calculation gives
\(d_5(R)>13/8\) and \(\chi(R)=1-d_5(R)/3<11/24\). The exact Haar
selection rule and the returned mean give the volume-independent row
estimate H11. Link-center symmetry makes the homogeneous connected heat
even. Cauchy\textquotesingle s formula therefore proves

\begingroup\small\begin{equation*}
\adjustbox{max width=\linewidth}{$\displaystyle \left\|\widehat C(\tau;\xi)-C_0(\tau)-\xi^2C_2(\tau)-\xi^4C_4(\tau)\right\|_{\mathrm{row}}
\leq\frac{67896}{169}e^{-13\tau/8}
\frac{(55|\xi|)^6}{1-(55|\xi|)^2},$}
\end{equation*}\endgroup

for every \(L\geq2\), \(\tau\geq0\) and \(|\xi|<1/55\), where
\(\|C\|_{\mathrm{row}}=\max_p\sum_q|C_{pq}|\). Column bounds agree by
transpose symmetry. The coefficients \(C_0,C_2,C_4\) are the unchanged
complete heat coefficients of the earlier edition; the improvement here
is the bound on their entire remaining tail.

The radius \(1/55\) exceeds both the preceding volume-heat radius
\(3/256\) and the saved fourth-reference source endpoint
\(0.018104972231644127076\). It is strictly inside \(\alpha_5\), so the
heat statement has its own open domain. For real coupling it requires
\(g^2>\sqrt{55}/2\), whereas the source and physical-gap construction
permits \(g^2\geq1/(2\sqrt{\alpha_5})\).

\subsection{Original moments and the complete physical
complement}\label{original-moments-and-the-complete-physical-complement}

Let \(R_{\mathrm{col}}c=\sum_pc_pr_p\) and
\(\Phi=\kappa A^{-1}R_{\mathrm{col}}\). On the centered physical space,

\begingroup\small\begin{equation*}
\adjustbox{max width=\linewidth}{$\displaystyle G_0=R_{\mathrm{col}}^*R_{\mathrm{col}},\quad
G_1=\kappa R_{\mathrm{col}}^*A^{-1}R_{\mathrm{col}},\quad
G_2=\Phi^*\Phi,\quad E=\Phi^*A\Phi=\kappa G_1.$}
\end{equation*}\endgroup

These distinguish forcing-state norm, inverse-energy response and the
state norm of the response. For \(k\geq1\), the spectral theorem gives

\begingroup\small\begin{equation*}
\adjustbox{max width=\linewidth}{$\displaystyle G_k=\frac1{(k-1)!}\int_0^\infty\tau^{k-1}\widehat C(\tau;\xi)\,d\tau.$}
\end{equation*}\endgroup

The leakage trial
\(J=(R_{\mathrm{col}}-3\Phi)^*(R_{\mathrm{col}}-3\Phi)=G_0-6G_1+9G_2\)
retains its signed time weight \(9\tau-6\). Integrating its absolute
value against \(e^{-13\tau/8}\) gives the improved tail prefactor
\(5156090136/3570125\) in H18. Together with the full inherited
degree-two and degree-four support sums, this supplies explicit widths
for the original matrices, uniformly at boundary rows as well as in the
interior.

The projection onto the response states and their coupling to the rest
of the physical space are

\begingroup\small\begin{equation*}
\adjustbox{max width=\linewidth}{$\displaystyle P=\Phi G_2^{-1}\Phi^*,\qquad Q=I-P,\qquad B_\Phi=QR_{\mathrm{col}}.$}
\end{equation*}\endgroup

For every centered physical form-domain vector \(\Phi c+h\),
\(h\in Q\mathcal H_0\), its full energy includes the mixed term:

\begingroup\small\begin{equation*}
\adjustbox{max width=\linewidth}{$\displaystyle q_A(\Phi c+h)=\kappa c^*G_1c
+2\kappa\operatorname{Re}\langle B_\Phi c,h\rangle+q_A(h).$}
\end{equation*}\endgroup

Let \(d_{\mathrm{ph}}>0\) be a stated rational lower bound for
\(d_5(\xi)\) at the coupling benchmark. H25 proves that \(D_Q=QAQ\) is
self-adjoint on \(\operatorname{Dom}(A)\cap Q\mathcal H_0\), with
\(D_Q\geq\kappa d_{\mathrm{ph}}I\). Minimizing over this entire
complement gives

\begingroup\small\begin{equation*}
\adjustbox{max width=\linewidth}{$\displaystyle \Phi_{\mathrm{full}}=\Phi-\kappa D_Q^{-1}B_\Phi,\qquad
E_{\mathrm{full}}=E-\kappa^2B_\Phi^*D_Q^{-1}B_\Phi,$}
\end{equation*}\endgroup

\begingroup\small\begin{equation*}
\adjustbox{max width=\linewidth}{$\displaystyle G_{\mathrm{full}}=G_2+\kappa^2B_\Phi^*D_Q^{-2}B_\Phi.$}
\end{equation*}\endgroup

Let \(P_R=R_{\mathrm{col}}G_0^{-1}R_{\mathrm{col}}^*\). The observed
state fraction is \(\|P_R\Phi c\|^2/\|\Phi c\|^2\) for \(c\ne0\). Two
improved benchmarks from §H5 are:

{\def\LTcaptype{none} % do not increment counter
\begin{longtable}[]{@{}
  >{\raggedright\arraybackslash}p{(\linewidth - 6\tabcolsep) * \real{0.2500}}
  >{\raggedleft\arraybackslash}p{(\linewidth - 6\tabcolsep) * \real{0.2500}}
  >{\raggedleft\arraybackslash}p{(\linewidth - 6\tabcolsep) * \real{0.2500}}
  >{\raggedleft\arraybackslash}p{(\linewidth - 6\tabcolsep) * \real{0.2500}}@{}}
\toprule\noalign{}
\begin{minipage}[b]{\linewidth}\raggedright
Original coupling
\end{minipage} & \begin{minipage}[b]{\linewidth}\raggedleft
Observed fraction exceeds
\end{minipage} & \begin{minipage}[b]{\linewidth}\raggedleft
Energy loss below
\end{minipage} & \begin{minipage}[b]{\linewidth}\raggedleft
State increase below
\end{minipage} \\
\midrule\noalign{}
\endhead
\bottomrule\noalign{}
\endlastfoot
\(g^2\geq13\) & \(999/1000\) & \(1/2000\) & \(1/2000\) \\
\(g^2\geq16\) & \(9999/10000\) & \(1/25000\) & \(1/25000\) \\
\end{longtable}
}

An energy loss below \(\eta\) means
\((1-\eta)E\prec E_{\mathrm{full}}\preceq E\); a state increase below
\(\eta\) means \(G_2\preceq G_{\mathrm{full}}\prec(1+\eta)G_2\). These
compare forms on the same coefficient space and are strict on nonzero
coefficients. They are not entrywise quotients. The table holds for all
\(L\geq2\) and \(a>0\). At \(g^2\geq16\), H29 also bounds the original
section-energy change by \(1/20000\).

\subsection{Spatial limits, finite heat time, and the continuum
boundary}\label{spatial-limits-finite-heat-time-and-the-continuum-boundary}

The original source and heat recurrences preserve connected finite
supports. When two boxes contain the same neighbourhood of the marked
faces, their corresponding low-degree coefficients coincide. The
remaining geometric tail gives full-sequence convergence of the local
vacuum means and heat correlations at fixed \(a,g\) with \(\xi<1/55\).
§H6 constructs the limiting physical state and self-adjoint generator
and passes the gap and bounded moment operators to the infinite spatial
lattice. At the benchmark couplings, the positive Gram lower bounds and
the complete complementary minimization pass to this limit as well.

For a finite time \(T\), the response \(\Phi_T=(I-e^{-TA})\Phi\) has the
explicit residual
\(A\Phi_T=\kappa R_{\mathrm{col}}-\kappa e^{-TA}R_{\mathrm{col}}\). If a
signed determinant return has total absolute rank \(\mathcal B\), H33
gives the physical choice

\begingroup\small\begin{equation*}
\adjustbox{max width=\linewidth}{$\displaystyle T=\frac1{\kappa d_{\mathrm{ph}}}\log\left(1+\frac{2\mathcal B}{\eta}\right)$}
\end{equation*}\endgroup

for error at most \(\eta>0\). The rank, physical time, mixed blocks and
quotient-minimum fibers remain explicit.

The simultaneous shrinking-spacing path in §H7 is \(a_n=a_0 2^{-n}\),
\(g_n^2=1/c_n\), \(c_n=g_0^{-2}+\beta n\log2\), so \(\xi_n=c_n^2/4\).
Source inclusion requires \(c_n\leq2\sqrt{\alpha_5}\); heat inclusion
requires \(c_n<2/\sqrt{55}\). For \(\beta>0\) that path eventually
leaves both domains. The fixed-spacing spatial limit therefore does not
prove a positive finite continuum mass or a nontrivial smooth
four-dimensional field. The sixth-source catalogue also remains absent
from the delivered source chain.

\subsection{Human sources and the two upstream
branches}\label{human-sources-and-the-two-upstream-branches}

The direct analytic antecedents used here are the
{e\allowbreak{}x\allowbreak{}p\allowbreak{}o\allowbreak{}n\allowbreak{}e\allowbreak{}n\allowbreak{}t\allowbreak{}i\allowbreak{}a\allowbreak{}l\allowbreak{}-\allowbreak{}v\allowbreak{}a\allowbreak{}c\allowbreak{}u\allowbreak{}u\allowbreak{}m\allowbreak{}/\allowbreak{}C\allowbreak{}a\allowbreak{}s\allowbreak{}i\allowbreak{}m\allowbreak{}i\allowbreak{}r\allowbreak{}}
construction {[}1{]}, the compact Fourier algebra {[}2{]}, and the
generator theorem {[}3{]}. The two manuscripts give their own
source-domain, endpoint and physical-return arguments. Their finite
exact checks support the declared algebraic identities and rational
inequalities; they do not constitute an external analytic review or a
formal verification of the infinite-dimensional passages.

Two separate earlier branches of the workbench credit Levent Alpöge. The
Jacobian branch begins with his original announcement of a
counterexample to global injectivity with constant nonzero Jacobian
{[}4{]}, which credits Akhil\textquotesingle s question and
Fable\textquotesingle s work. Tao\textquotesingle s later exposition
{[}5{]} is a separate source. The \(S^6\) branch begins with the claimed
complex-threefold construction circulated by Alpöge with AI assistance
{[}6{]}; Engel\textquotesingle s exposition {[}7{]} separately credits
that source. These constructions and their later coordinate calculations
are distinct. Neither originating construction is invoked as a direct
lemma in the fifth-source or heat-complement proofs reproduced here. The
workbench\textquotesingle s
\href{https://github.com/KokunoYumeto/yang-mills-interacting-workbench/blob/main/ATTRIBUTION.md}{attribution
record} records the earlier points of use.

\subsection{References}\label{references}

\begin{enumerate}
\def\labelenumi{\arabic{enumi}.}
\tightlist
\item
  D. Schütte, Zheng Weihong and C. J. Hamer, \emph{The Coupled Cluster
  Method in Hamiltonian Lattice Field Theory}, 1996, Sections 2-5.
  \href{https://arxiv.org/abs/hep-lat/9603026v1}{{a\allowbreak{}r\allowbreak{}X\allowbreak{}i\allowbreak{}v\allowbreak{}:\allowbreak{}h\allowbreak{}e\allowbreak{}p\allowbreak{}-\allowbreak{}l\allowbreak{}a\allowbreak{}t\allowbreak{}/\allowbreak{}9\allowbreak{}6\allowbreak{}0\allowbreak{}3\allowbreak{}0\allowbreak{}2\allowbreak{}6\allowbreak{}v\allowbreak{}1\allowbreak{}}}.
\item
  P. Eymard, \emph{L\textquotesingle algèbre de Fourier
  d\textquotesingle un groupe localement compact}, Bulletin de la
  Société Mathématique de France 92 (1964), 181-236.
  \href{https://www.numdam.org/item/BSMF_1964__92__181_0/}{Original
  article}.
\item
  G. Lumer and R. S. Phillips, \emph{Dissipative operators in a Banach
  space}, Pacific Journal of Mathematics 11 (1961), 679-698, Theorem
  3.1.
  \href{https://msp.org/pjm/1961/11-2/pjm-v11-n2-p19-s.pdf}{Original
  article}.
\item
  Levent Alpöge,
  \href{https://x.com/__alpoge__/status/2079028340955197566}{original
  Jacobian counterexample announcement}, 20 July 2026; credits Akhil and
  Fable.
\item
  Terence Tao, \emph{A digestion of the Jacobian conjecture
  counterexample}, 21 July 2026.
  \href{https://terrytao.wordpress.com/2026/07/21/a-digestion-of-the-jacobian-conjecture-counterexample/}{Exposition}.
\item
  Levent Alpöge, \(S^6\) construction circulated with AI assistance.
  \href{https://alpo.ge/s6.pdf}{Original manuscript}, especially p. 2,
  Setup.
\item
  Philip Engel, \emph{Complex structures on \(S^6\)}, 13 September 2026.
  \href{https://philip-engel.github.io/S6.pdf}{Expository manuscript},
  abstract and Section 1.3.
\end{enumerate}

\clearpage
\appendix

\section{The complete fifth reference: original spin channels, a
convergent correction, and physical
energy}\label{the-complete-fifth-reference-original-spin-channels-a-convergent-correction-and-physical-energy}

21 September 2026. This continues the actual Yang--Mills checkpoint in
the cumulative volume/heat edition. The fifth coefficient catalogue is
an input already calculated there. Here its complete signed channel
action is evaluated for quantitative source control, then the entire
remaining nonlinear equation is solved on an explicit disk. The Riemann
attachments have no mathematical role in this calculation. The proof
supplies written analytic arguments and exact finite certificates; it
does not assert an independent external analytical review, a
formal-assistant proof, historical priority, or a continuum mass gap.

\subsection{F1. Original coordinates, spaces, and coefficient
norms}\label{f1.-original-coordinates-spaces-and-coefficient-norms}

Fix an integer L\textgreater=2. The original vertices are
\{-L,...,L\}\^{}3, with every contained positively oriented
nearest-neighbor edge e=(n,i) and every contained elementary face
p=(n;i,j), i\textless j. Keep

\begingroup\small\begin{equation*}
\adjustbox{max width=0.92\linewidth}{$\displaystyle W_p=\operatorname{tr}\{U_i(n)U_j(n+e_i)U_i(n+e_j)^{-1}U_j(n)^{-1}\},\quad
T_a=-i\sigma_a/2,\quad X_{e,a}f=\left.\partial_t f(\ldots,e^{tT_a}U_e,\ldots)\right|_0.$}\tag{F1}
\end{equation*}\endgroup

The physical Hilbert space is the invariant part of original
product-Haar-probability L2, with all vertex transformations, including
boundary vertices, included. The original operator is

\begingroup\small\begin{equation*}
\adjustbox{max width=0.92\linewidth}{$\displaystyle H_L=\kappa K+\kappa\xi(2M-S),\quad K=-\sum_{e,a}X_{e,a}^2,
\quad S=\sum_pW_p,\quad M=|\mathsf P_L|,
\quad\kappa=\frac{2g^2}{a},\quad\xi=\frac1{4g^4},\quad a,g>0.$}\tag{F2}
\end{equation*}\endgroup

The physical operator domain is H2 and the form domain H1 on this
compact product. Its smooth bounded potential preserves self-adjointness
and compact resolvent. Write P\_H f=int f dU, Q\_H=I-P\_H. For the
actual positive ground vector psi,

\begingroup\small\begin{equation*}
\adjustbox{max width=0.92\linewidth}{$\displaystyle q_{H-E_0}(\psi f)=\kappa\int\psi^2\sum_i|X_if|^2dU.$}\tag{F3}
\end{equation*}\endgroup

Integration by parts using H psi=E0 psi proves this identity. The
modulus inequality, elliptic regularity and the maximum principle give a
positive ground vector; (F3) gives uniqueness. Gauge transformations
consequently fix this positive unit vector.

For the original product representation pi\_j, with
{d\allowbreak{}\_\allowbreak{}j\allowbreak{}=\allowbreak{}p\allowbreak{}r\allowbreak{}o\allowbreak{}d\allowbreak{}u\allowbreak{}c\allowbreak{}t\allowbreak{}\_\allowbreak{}e\allowbreak{}(\allowbreak{}2\allowbreak{}j\allowbreak{}\_\allowbreak{}e\allowbreak{}+\allowbreak{}1\allowbreak{})\allowbreak{},\allowbreak{}}
retain

\begingroup\small\begin{equation*}
\adjustbox{max width=0.92\linewidth}{$\displaystyle A_j(f)=d_j\int f(U)\pi_j(U)^*dU,\quad
f=\sum_j\operatorname{Tr}(A_j(f)\pi_j),\quad
c_j=\sum_e j_e(j_e+1),\quad \|f\|_X=\sum_j\|A_j(f)\|_1.$}\tag{F4}
\end{equation*}\endgroup

The norm on the right is an auxiliary coefficient estimate; the physical
pairing remains (F3). Original Haar matrix-coefficient orthogonality
proves the Fourier map and inverse. Decomposing a product representation
with all multiplicities, pinching its coefficient matrix, and tracing
its multiplicity spaces gives

\begingroup\small\begin{equation*}
\adjustbox{max width=0.92\linewidth}{$\displaystyle \|fh\|_X\le\|f\|_X\|h\|_X,
\qquad \|A_j(X_{e,a}f_j)\|_1\le j_e\|A_j(f_j)\|_1.$}\tag{F5}
\end{equation*}\endgroup

Pinching is an average of unitary conjugations. For the partial trace,
duality with Z tensor I, \textbar\textbar Z\textbar\textbar\textless=1,
proves contraction of trace norm. This proves the needed compact-group
Fourier-algebra estimates directly; Eymard\textquotesingle s Fourier
algebra is the antecedent.

Each original coefficient also retains a finite edge label S containing
its active representation support. Set

\begingroup\small\begin{equation*}
\adjustbox{max width=0.92\linewidth}{$\displaystyle \begin{aligned}
\|v\|_{\rm loc}&=\sup_a\sum_{S\ni a,j\ne0}c_j\|A_{S,j}\|_1,\\
m(v)&=\sup_a\sum_{S\ni a,j\ne0}\Big(\sum_ej_e\Big)\|A_{S,j}\|_1,\\
t(v)&=\sup_e\sum_{S\ni e,j\ne0}j_e\|A_{S,j}\|_1.
\end{aligned}$}\tag{F6}
\end{equation*}\endgroup

Assembly sums A\_(S,j) over all labels containing supp(j). Its kernel
consists exactly of the equations that these sums vanish. For each
enlarged S, move its actual A\_(S,j) to supp(j) and retain the two-entry
relation (+A at S,-A at supp(j)). Reading the enlarged-label entries is
the inverse decomposition on this kernel. Hence enlarged zero-function
labels remain represented.

\subsection{F2. Gauge estimate and bilinear
source}\label{f2.-gauge-estimate-and-bilinear-source}

At each original vertex, a nonzero physical coefficient contains a
singlet intertwiner. Grouping the j\_e factor gives
j\_e\textless=sum\_(f incident v,f!=e)j\_f by the highest weight of the
remaining tensor factors. For a fixed edge e=\{u,v\}, the other edges at
u,v have total spin J1\textgreater=2j\_e. Their outer endpoints are
distinct: a common endpoint would form a triangle in the original cubic
graph. Summing their own singlet inequalities gives J1\textless=2J2,
counting every further edge at most twice. Thus sum\_f
{j\allowbreak{}\_\allowbreak{}f\allowbreak{}\textgreater{}\allowbreak{}=\allowbreak{}j\allowbreak{}\_\allowbreak{}e\allowbreak{}+\allowbreak{}J\allowbreak{}1\allowbreak{}+\allowbreak{}J\allowbreak{}2\allowbreak{}\textgreater{}\allowbreak{}=\allowbreak{}4\allowbreak{}j\allowbreak{}\_\allowbreak{}e\allowbreak{}.\allowbreak{}}
Since j(j+1)\textgreater=(3/2)j at every nonzero half-integer,

\begingroup\small\begin{equation*}
\adjustbox{max width=0.92\linewidth}{$\displaystyle c_j\ge6j_e\quad\hbox{on physical blocks},\qquad
\sum_e j_e\le\frac23c_j,\qquad
m(v)\le\frac23\|v\|_{\rm loc},\quad t(v)\le\frac16\|v\|_{\rm loc}.$}\tag{F7}
\end{equation*}\endgroup

This proves c\_j\textgreater=3 on the nonconstant physical space. The
fundamental elementary plaquette attains 3. Active bridges and open
boundaries remain in this argument.

Define the original bilinear differential and its zero-Haar inverse,

\begingroup\small\begin{equation*}
\adjustbox{max width=0.92\linewidth}{$\displaystyle \Gamma(f,h)=\sum_{e,a}(X_{e,a}f)(X_{e,a}h),\qquad
B(f,h)=K^{-1}Q_H\Gamma(f,h).$}\tag{F8}
\end{equation*}\endgroup

Every input pair retains its union label. Every removed Haar scalar is
stored as a scalar energy contribution. Bounding anchors in each input
separately, (F5) gives

\begingroup\small\begin{equation*}
\adjustbox{max width=0.92\linewidth}{$\displaystyle \|B(f,h)\|_{\rm loc}\le3\{m(f)t(h)+m(h)t(f)\}
\le\frac23\|f\|_{\rm loc}\|h\|_{\rm loc}.$}\tag{F9}
\end{equation*}\endgroup

The factor three is the complete generator sum. The output c\_j cancels
exactly against K\^{}\{-1\} on the nonconstant output. Two useful
consequences are

\begingroup\small\begin{equation*}
\adjustbox{max width=0.92\linewidth}{$\displaystyle \|2B(q,\cdot)\|_{\rm loc\to loc}\le m(q)+4t(q),\qquad
\|2\Gamma(v,h)\|_X\le4t(v)\|Kh\|_X.$}\tag{F10}
\end{equation*}\endgroup

The latter includes the full scalar output.

\subsection{F3. Orientation-resolved trace coefficient and its exact
transport}\label{f3.-orientation-resolved-trace-coefficient-and-its-exact-transport}

Here is a new finite coefficient bound used before absolute summation.
First regard each of the l occurrences in a trace word as an independent
d-dimensional representation variable, with orientation signs s\_i=+1 or
-1. Replace each inverse by its actual invariant-dual transpose C
pi(U)\^{}T C\^{}\{-1\}, where C is unitary (for fundamental SU(2), C is
the alternating two-index matrix). These coefficient-side unitary
matrices preserve the singular values. In the remaining integer index
tensor, the cyclic trace indices r\_i contribute

\begingroup\small\begin{equation*}
\adjustbox{max width=0.92\linewidth}{$\displaystyle A_{(b_i),(a_i)}=\sum_{r_1,\ldots,r_l}
\prod_i\mathbf1\{(a_i,b_i)=(r_i,r_{i+1})\text{ for }s_i=+1;
\ (a_i,b_i)=(r_{i+1},r_i)\text{ for }s_i=-1\},\quad r_{l+1}=r_1.$}\tag{F11}
\end{equation*}\endgroup

Let 2k be the number of cyclic sign changes. At a vertex with equal
consecutive signs, r\_i appears once in a row and once in a column; this
gives an identity factor. A + to - transition gives a two-row delta
vector; a - to + transition gives a two-column delta vector. There are k
of each, each of norm sqrt(d), and l-2k identity factors. Separate row
and column permutations consequently give

\begingroup\small\begin{equation*}
\adjustbox{max width=0.92\linewidth}{$\displaystyle A^*A\text{ has nonzero eigenvalue }d^{2k}
\text{ of multiplicity }d^{l-2k},\qquad
\boxed{\|A\|_1=d^{l-k}.}$}\tag{F12}
\end{equation*}\endgroup

Equivalently (A* A)\^{}2=d\^{}(2k)A* A and Tr(A* A)=d\^{}l. This
includes every original color sum. The checker constructs the actual
integer matrices for all signs through length seven at d=2 and length
five at d=3: 316 cases, with both complete matrix identities checked.

Repeated occurrences of a link are now identified by the diagonal
compact-group homomorphism sending the one original U\_e to its
occurrence tuple. The representation on that tuple decomposes unitarily
with all multiplicities. The pinching and partial-trace argument of (F5)
proves that this actual pullback contracts the Fourier trace norm. Thus
2\^{}(l-k) is a rigorous upper bound for an original fundamental trace
word, and products multiply these bounds.

The original lattice coordinate transport is y\_i=epsilon\_i
x\_(pi(i))-d\_i, with inverse
{x\allowbreak{}\_\allowbreak{}(\allowbreak{}p\allowbreak{}i\allowbreak{}(\allowbreak{}i\allowbreak{})\allowbreak{})\allowbreak{}=\allowbreak{}e\allowbreak{}p\allowbreak{}s\allowbreak{}i\allowbreak{}l\allowbreak{}o\allowbreak{}n\allowbreak{}\_\allowbreak{}i\allowbreak{}(\allowbreak{}y\allowbreak{}\_\allowbreak{}i\allowbreak{}+\allowbreak{}d\allowbreak{}\_\allowbreak{}i\allowbreak{})\allowbreak{}.\allowbreak{}}
It sends the original U\_e to the mapped U or its inverse according to
orientation. The underlying physical Haar norm and Casimir intertwine
through that map. For the coefficient estimate, instead of assigning
invariance to every partial transpose, we explicitly retain all eight
sign reflections. For an original trace polynomial P=sum\_m c\_m
product\_(w in m)Tr(w), put

\begingroup\small\begin{equation*}
\adjustbox{max width=0.92\linewidth}{$\displaystyle \mathfrak b(P)=
\max_{\epsilon\in\{\pm1\}^3}
\sum_m|c_m|\,2^{\sum_{w\in m}(|w|-k_\epsilon(w))}.$}\tag{F13}
\end{equation*}\endgroup

Here k\_epsilon(w) counts half the cyclic sign changes after the
original axis reflection. The axes permutation only relabels the three
reflected coordinates. This proves a bound simultaneously on every
original signed transport. Every nonempty closed cubic word has both
signs after any reflection, since its net displacement in each
coordinate is zero. Hence k\_epsilon\textgreater=1 and (F13) is at most
the preceding length-only sum
{s\allowbreak{}u\allowbreak{}m\allowbreak{}\_\allowbreak{}m\allowbreak{}|\allowbreak{}c\allowbreak{}\_\allowbreak{}m\allowbreak{}|\allowbreak{}2\allowbreak{}\textasciicircum{}\allowbreak{}(\allowbreak{}t\allowbreak{}o\allowbreak{}t\allowbreak{}a\allowbreak{}l\allowbreak{}}
length-number of traces).

For clarity, on four independent factors, partial transpose of the last
two tensor coordinates sends the all-positive coefficient matrix of
trace norm 16 to the ++-\/- matrix of trace norm 8. The executable
checks this exact map and both costs. No norm identity under that
partial transpose is inserted.

\subsection{F4. Full fifth-source inputs and signed spin-channel
recurrence}\label{f4.-full-fifth-source-inputs-and-signed-spin-channel-recurrence}

The preserved fifth catalogue contains 662 representatives, 124864
anchored original multisets and 14063 rational trace-monomial terms.
Every coefficient v\_nu has its complete original edges, words and
multiplicities. The original coefficient equation is

\begingroup\small\begin{equation*}
\adjustbox{max width=0.92\linewidth}{$\displaystyle v_1=\frac13\sum_pW_p,\quad
v_n=\sum_{i=1}^{n-1}B(v_i,v_{n-i}),\qquad
v_5=2B(v_1,v_4)+2B(v_2,v_3).$}\tag{F14}
\end{equation*}\endgroup

For 321 of the fifth representatives all five faces are distinct and
each original edge occurs at most twice. Let E\_s be the repeated edges;
each carries final spin b\_e=0 or 1. A singly occurring edge carries
1/2. For a subset A of those five faces, let r\_e(A) count its original
occurrences and define

\begingroup\small\begin{equation*}
\adjustbox{max width=0.92\linewidth}{$\displaystyle c_b(A)=\sum_{r_e(A)=1}\frac34+
\sum_{r_e(A)=2}b_e(b_e+1).$}\tag{F15}
\end{equation*}\endgroup

The complete admissible spin assignments retain every original vertex
singlet constraint. For c\_b(A)\textgreater0 define

\begingroup\small\begin{equation*}
\adjustbox{max width=0.92\linewidth}{$\displaystyle a_b(\{p\})=\frac13,\qquad
a_b(A)=\frac1{2c_b(A)}
\sum_{\varnothing\ne B\subsetneq A}
[c_b(B)+c_b(A\setminus B)-c_b(A)]a_b(B)a_b(A\setminus B).$}\tag{F16}
\end{equation*}\endgroup

The empty coefficient and zero-Casimir coefficient are zero in the
zero-Haar source. For these distinct-face proper subsets there is no
nonempty closed surface. Each intermediate shared edge has either the
same final pair of fundamental occurrences, or one fundamental
occurrence. Therefore its pair projector and its actual intermediate
Casimir intertwine with the final projector. Projectors on different
edges commute. Applying

\begingroup\small\begin{equation*}
\adjustbox{max width=0.92\linewidth}{$\displaystyle 2\Gamma(f,h)=(c_f+c_h)fh-K(fh)$}\tag{F17}
\end{equation*}\endgroup

to every ordered input split proves

\begingroup\small\begin{equation*}
\adjustbox{max width=0.92\linewidth}{$\displaystyle \boxed{v_\nu=\sum_b a_b(A_{\rm all})P_b\prod_{p\in A_{\rm all}}W_p.}$}\tag{F18}
\end{equation*}\endgroup

Every scalar in (F16) is evaluated as a rational number before taking
its absolute value. The formula retains all input orderings and zero
channels. Complete original quaternion polynomial identities verify
(F18) against the inherited fifth coefficient for all 321
representatives, rather than testing only values at selected
configurations.

For a set Z of repeated edges,

\begingroup\small\begin{equation*}
\adjustbox{max width=0.92\linewidth}{$\displaystyle P_{0,Z}=\prod_{e\in Z}(I-E_e/2),\qquad E_e=-\sum_aX_{e,a}^2.$}\tag{F19}
\end{equation*}\endgroup

This is the original spin-zero projector on the degree-two occurrence
space. Let x\_b be the nonnegative trace norm of the coefficient block
of P\_b product W\_p. Pinching into the complete edge-spin blocks gives

\begingroup\small\begin{equation*}
\adjustbox{max width=0.92\linewidth}{$\displaystyle \sum_{b:\ b_e=0\ (e\in Z)}x_b
=\|P_{0,Z}\prod_pW_p\|_X
\le\mathfrak b(P_{0,Z}\prod_pW_p).$}\tag{F20}
\end{equation*}\endgroup

Each right side is calculated from the complete projected trace
polynomial and all eight reflection bounds. The 321 families contain
6240 such constraints; 3498 are strictly tighter than their length-only
predecessors.

The three objective families are

\begingroup\small\begin{equation*}
\adjustbox{max width=0.92\linewidth}{$\displaystyle \sum_b|a_b|c_bx_b,\quad
\sum_b|a_b|(\sum_ej_e(b))x_b,\quad
\sum_b|a_b|j_e(b)x_b\quad\text{for every original edge }e.$}\tag{F21}
\end{equation*}\endgroup

Each supplied rational dual lambda\textgreater=0 satisfies A\^{}T
lambda\textgreater=the actual objective and has cost b\^{}T lambda.
Hence its cost bounds that objective for the actual x. The records also
give feasible primal vectors with the same value. A separate auditor
reconstructs the constraints, objectives, original Casimirs and
polynomial projections, then verifies all 5726 primal/dual certificates
without invoking the optimization solver. They certify bounds on the
actual coefficients, not an assumption that the actual coefficient
masses attain the finite polytope maximum.

\subsection{F5. The remaining original configurations and complete
anchored
return}\label{f5.-the-remaining-original-configurations-and-complete-anchored-return}

For every ordered split of a fifth multiset into original submultisets
A,B, the complete cubic channel polytopes and the exact saved fourth
coefficients give

\begingroup\small\begin{equation*}
\adjustbox{max width=0.92\linewidth}{$\displaystyle \|B(v_A,v_B)\|_{\rm loc,one\ label}
\le3\sum_e t_e(v_A)t_e(v_B).$}\tag{F22}
\end{equation*}\endgroup

For a (2,3) split the record also maximizes the full bilinear expression
over the two actual finite input-channel polytopes. Its variables,
original edge spins, signed scalar coefficients and all polytope
vertices are retained in the inherited channel source. For a (1,4)
split, t\_e(v\_1)=4/3 on that face; hence the sharper joint bound is

\begingroup\small\begin{equation*}
\adjustbox{max width=0.92\linewidth}{$\displaystyle 4\sup\left\{\sum_j\sum_{e\in\partial p\cap\operatorname{supp}B}
 j_e\|A_j(v_B)\|_1\right\}.$}\tag{F23}
\end{equation*}\endgroup

This simultaneous edge objective is evaluated on the actual fourth
channel polytope where available. On every fourth row it is bounded as
well by its saved local-Casimir bound times
{m\allowbreak{}a\allowbreak{}x\allowbreak{}\_\allowbreak{}j\allowbreak{}(\allowbreak{}s\allowbreak{}u\allowbreak{}m\allowbreak{}\_\allowbreak{}s\allowbreak{}e\allowbreak{}l\allowbreak{}e\allowbreak{}c\allowbreak{}t\allowbreak{}e\allowbreak{}d\allowbreak{}}
j\_e/c\_j), by the complete raw trace estimate, and by the sum of its
saved individual edge bounds. For a repeated single face its complete
spin-one and spin-two characters give the additional exact
coefficient-norm bound. Each of these is a bound on the same scalar
objective, so the minimum of the scalar upper bounds is valid.

Summing every ordered split proves an upper bound C\_nu for the actual
fifth coefficient on its original union. Its total and marked-spin
bounds also follow from the full admissible output-spin list:

\begingroup\small\begin{equation*}
\adjustbox{max width=0.92\linewidth}{$\displaystyle m_\nu\le C_\nu\max_j\frac{\sum_ej_e}{c_j},\qquad
t_{\nu,e}\le C_\nu\max_j\frac{j_e}{c_j}.$}\tag{F24}
\end{equation*}\endgroup

The complete inherited trace expression supplies independent raw upper
bounds; the explicit one-face fifth character supplies its own further
bound. On the 321 rows of F4, the complete signed channel objectives
(F21) are also used. Thus all 662 rows are covered, including the other
341 original
{m\allowbreak{}u\allowbreak{}l\allowbreak{}t\allowbreak{}i\allowbreak{}p\allowbreak{}l\allowbreak{}i\allowbreak{}c\allowbreak{}i\allowbreak{}t\allowbreak{}y\allowbreak{}/\allowbreak{}s\allowbreak{}h\allowbreak{}a\allowbreak{}r\allowbreak{}e\allowbreak{}d\allowbreak{}-\allowbreak{}e\allowbreak{}d\allowbreak{}g\allowbreak{}e\allowbreak{}}
configurations. No generic distinct-face recurrence is applied to an
edge with three or more occurrences.

The original anchor is the edge from (0,0,0) to (1,0,0). For every one
of the 124864 original anchored multisets the saved signed-coordinate
map sends both endpoints into the representative edge set. Summing that
representative\textquotesingle s actual marked-edge bound, rather than a
freely selected representative edge, proves

\begingroup\small\begin{equation*}
\adjustbox{max width=0.92\linewidth}{$\displaystyle \begin{aligned}
\|v_5\|_{\rm loc}&\le
\frac{1141532378446937587064431}{460877915422606500}<2476866,\\
m(v_5)&\le
\frac{6682496758630785752687413327666759}{4077966203261057899505463750}<1638684,\\
t(v_5)&\le
\frac{126023896114060026308152742875366210205679497747117260778732756243249}
{662840019461690772503734890389307666881000651173392960000000000}<190128.
\end{aligned}$}\tag{F25}
\end{equation*}\endgroup

Every boundary box uses a subset of these original labels, so the same
bounds hold independently of L. The complete per-row and per-anchor sums
are in generated/rows and
{g\allowbreak{}e\allowbreak{}n\allowbreak{}e\allowbreak{}r\allowbreak{}a\allowbreak{}t\allowbreak{}e\allowbreak{}d\allowbreak{}/\allowbreak{}f\allowbreak{}i\allowbreak{}f\allowbreak{}t\allowbreak{}h\allowbreak{}\_\allowbreak{}b\allowbreak{}u\allowbreak{}d\allowbreak{}g\allowbreak{}e\allowbreak{}t\allowbreak{}s\allowbreak{}.\allowbreak{}j\allowbreak{}s\allowbreak{}o\allowbreak{}n\allowbreak{}.\allowbreak{}}
The 124864 coordinate transports and full fifth-source polynomials are
preserved in the preceding directory.

\subsection{F6. Full fifth-reference residual, inverse, and nonlinear
correction}\label{f6.-full-fifth-reference-residual-inverse-and-nonlinear-correction}

For the following closed formulas use the proved upward bounds

\begingroup\small\begin{equation*}
\adjustbox{max width=0.92\linewidth}{$\displaystyle \begin{aligned}
(m_1,t_1)&=(64/3,16/3),\\
(m_2,t_2)&=(5834/39,137/6),\\
(m_3,t_3)&=(336572872/208845,225985217/1253070),\\
(m_4,t_4)&=(17270702970768271/341697152160,
110695177857394584026401/18025447358750832000),\\
(m_5,t_5)&=(1638684,190128).
\end{aligned}$}\tag{F26}
\end{equation*}\endgroup

Each coefficient is a bound on the same original source norms in F6. Set
q5=sum\_(i=1)\^{}5 xi\^{}i v\_i. Its complete residual is

\begingroup\small\begin{equation*}
\adjustbox{max width=0.92\linewidth}{$\displaystyle \begin{aligned}
R_5={}&\xi v_1+B(q_5,q_5)-q_5\\
={}&\xi^6[2B(v_1,v_5)+2B(v_2,v_4)+B(v_3,v_3)]\\
&+\xi^7[2B(v_2,v_5)+2B(v_3,v_4)]\\
&+\xi^8[2B(v_3,v_5)+B(v_4,v_4)]
+2\xi^9B(v_4,v_5)+\xi^{10}B(v_5,v_5).
\end{aligned}$}\tag{F27}
\end{equation*}\endgroup

No coefficient above degree five is treated as zero. Define

\begingroup\small\begin{equation*}
\adjustbox{max width=0.92\linewidth}{$\displaystyle b_{ij}=3(m_it_j+m_jt_i),\quad
\ell(x)=\sum_{i=1}^5(m_i+4t_i)x^i,\quad
\delta(x)=\sum_{\substack{1\le i,j\le5\\i+j\ge6}}b_{ij}x^{i+j},\quad
\mathscr D(x)=(1-\ell(x))^2-\frac83\delta(x).$}\tag{F28}
\end{equation*}\endgroup

All coefficients of ell and delta are positive. On {[}0,19/1000{]},
ell\textless1, so D\textquotesingle=
{-\allowbreak{}2\allowbreak{}(\allowbreak{}1\allowbreak{}-\allowbreak{}e\allowbreak{}l\allowbreak{}l\allowbreak{})\allowbreak{}e\allowbreak{}l\allowbreak{}l\allowbreak{}'\allowbreak{}-\allowbreak{}(\allowbreak{}8\allowbreak{}/\allowbreak{}3\allowbreak{})\allowbreak{}d\allowbreak{}e\allowbreak{}l\allowbreak{}t\allowbreak{}a\allowbreak{}'\allowbreak{}\textless{}\allowbreak{}0\allowbreak{}.\allowbreak{}}
The signs at 18/1000 and 19/1000 are opposite. Therefore there is
exactly one root alpha in that interval and it is the first positive
root. Exact rational bisection and integer-square bounds give

\begingroup\small\begin{equation*}
\adjustbox{max width=0.92\linewidth}{$\displaystyle \begin{aligned}
0.018424953576117616681&<\alpha<0.018424953576117616682,\\
3.683551983985727304439&<\frac1{2\sqrt\alpha}
<3.683551983985727304440.
\end{aligned}$}\tag{F29}
\end{equation*}\endgroup

For any complex \textbar xi\textbar\textless=alpha, J5=2B(q5,.)
satisfies \textbar\textbar J5\textbar\textbar\textless=ell\textless1.
Thus L5=I-J5 has the actual bounded inverse
sum\_(n\textgreater=0)J5\^{}n, on the original local labelled source.
The full correction w=v-q5 solves

\begingroup\small\begin{equation*}
\adjustbox{max width=0.92\linewidth}{$\displaystyle L_5w=R_5+B(w,w).$}\tag{F30}
\end{equation*}\endgroup

Put w1=L5\^{}\{-1\}R5 and
{w\allowbreak{}\_\allowbreak{}n\allowbreak{}=\allowbreak{}L\allowbreak{}5\allowbreak{}\textasciicircum{}\allowbreak{}\{\allowbreak{}-\allowbreak{}1\allowbreak{}\}\allowbreak{}s\allowbreak{}u\allowbreak{}m\allowbreak{}\_\allowbreak{}(\allowbreak{}i\allowbreak{}+\allowbreak{}j\allowbreak{}=\allowbreak{}n\allowbreak{})\allowbreak{}B\allowbreak{}(\allowbreak{}w\allowbreak{}\_\allowbreak{}i\allowbreak{},\allowbreak{}w\allowbreak{}\_\allowbreak{}j\allowbreak{})\allowbreak{}.\allowbreak{}}
Its nth norm is bounded by
{C\allowbreak{}a\allowbreak{}t\allowbreak{}\_\allowbreak{}(\allowbreak{}n\allowbreak{}-\allowbreak{}1\allowbreak{})\allowbreak{}[\allowbreak{}(\allowbreak{}2\allowbreak{}/\allowbreak{}3\allowbreak{})\allowbreak{}/\allowbreak{}(\allowbreak{}1\allowbreak{}-\allowbreak{}e\allowbreak{}l\allowbreak{}l\allowbreak{})\allowbreak{}]\allowbreak{}\textasciicircum{}\allowbreak{}(\allowbreak{}n\allowbreak{}-\allowbreak{}1\allowbreak{})\allowbreak{}[\allowbreak{}d\allowbreak{}e\allowbreak{}l\allowbreak{}t\allowbreak{}a\allowbreak{}/\allowbreak{}(\allowbreak{}1\allowbreak{}-\allowbreak{}e\allowbreak{}l\allowbreak{}l\allowbreak{})\allowbreak{}]\allowbreak{}\textasciicircum{}\allowbreak{}n\allowbreak{}.\allowbreak{}}
This complete series converges even at the endpoint, and

\begingroup\small\begin{equation*}
\adjustbox{max width=0.92\linewidth}{$\displaystyle \boxed{\|w\|_{\rm loc}\le w_*(x)
=\frac34[1-\ell(x)-\sqrt{\mathscr D(x)}],\qquad x=|\xi|\le\alpha.}$}\tag{F31}
\end{equation*}\endgroup

With
{t\allowbreak{}h\allowbreak{}e\allowbreak{}t\allowbreak{}a\allowbreak{}=\allowbreak{}8\allowbreak{}d\allowbreak{}e\allowbreak{}l\allowbreak{}t\allowbreak{}a\allowbreak{}/\allowbreak{}[\allowbreak{}3\allowbreak{}(\allowbreak{}1\allowbreak{}-\allowbreak{}e\allowbreak{}l\allowbreak{}l\allowbreak{})\allowbreak{}\textasciicircum{}\allowbreak{}2\allowbreak{}]\allowbreak{}\textless{}\allowbreak{}=\allowbreak{}1\allowbreak{},\allowbreak{}}
its omitted tree orders satisfy

\begingroup\small\begin{equation*}
\adjustbox{max width=0.92\linewidth}{$\displaystyle \left\|\sum_{n>N}w_n\right\|_{\rm loc}
\le\frac34(1-\ell)\theta^{N+1}\frac{\binom{2N}{N}}{4^N}
\le\frac34(1-\ell)\frac{\theta^{N+1}}{\sqrt{N+1}}.$}\tag{F32}
\end{equation*}\endgroup

The identity
{C\allowbreak{}a\allowbreak{}t\allowbreak{}\_\allowbreak{}N\allowbreak{}/\allowbreak{}4\allowbreak{}\textasciicircum{}\allowbreak{}N\allowbreak{}=\allowbreak{}2\allowbreak{}[\allowbreak{}b\allowbreak{}i\allowbreak{}n\allowbreak{}o\allowbreak{}m\allowbreak{}(\allowbreak{}2\allowbreak{}N\allowbreak{},\allowbreak{}N\allowbreak{})\allowbreak{}/\allowbreak{}4\allowbreak{}\textasciicircum{}\allowbreak{}N\allowbreak{}-\allowbreak{}b\allowbreak{}i\allowbreak{}n\allowbreak{}o\allowbreak{}m\allowbreak{}(\allowbreak{}2\allowbreak{}N\allowbreak{}+\allowbreak{}2\allowbreak{},\allowbreak{}N\allowbreak{}+\allowbreak{}1\allowbreak{})\allowbreak{}/\allowbreak{}4\allowbreak{}\textasciicircum{}\allowbreak{}(\allowbreak{}N\allowbreak{}+\allowbreak{}1\allowbreak{})\allowbreak{}]\allowbreak{}}
proves the tail exactly; the central binomial bound follows by induction
after squaring. Thus the endpoint uses an explicit convergent tail
rather than a strict endpoint contraction.

On every finite graph, the total c-weighted coefficient sum is at most
\textbar E\_L\textbar{} times the local norm. First and second
derivatives converge because j\_e and j\_e j\_f are bounded by constant
multiples of c\_j. The assembled source therefore solves

\begingroup\small\begin{equation*}
\adjustbox{max width=0.92\linewidth}{$\displaystyle Kv=\xi S+\Gamma(v,v)-P_H\Gamma(v,v).$}\tag{F33}
\end{equation*}\endgroup

For real xi, exp(v) is positive and gives

\begingroup\small\begin{equation*}
\adjustbox{max width=0.92\linewidth}{$\displaystyle \psi=e^{v+c_L},\quad c_L=-\frac12\log\int e^{2v}dU,
\quad E_0=2\kappa M\xi-\kappa P_H\Gamma(v,v).$}\tag{F34}
\end{equation*}\endgroup

Substitution into the original H gives its exact eigen-equation.
Equation F3 proves this eigenvalue is the bottom, and uniqueness
identifies the actual vacuum. Elliptic bootstrapping gives all higher
derivatives. For complex xi the same nonzero exponential solves the
complex equation; the mean functional and its nonvanishing denominator
are proved in the heat companion.

\subsection{F7. Actual physical spectral and primitive
return}\label{f7.-actual-physical-spectral-and-primitive-return}

Let t\_* =sum\_i t\_i x\^{}i+w\_*/6 and chi=4t\_*. Equations F10 and F31
give
{|\allowbreak{}|\allowbreak{}2\allowbreak{}G\allowbreak{}a\allowbreak{}m\allowbreak{}m\allowbreak{}a\allowbreak{}(\allowbreak{}v\allowbreak{},\allowbreak{}h\allowbreak{})\allowbreak{}|\allowbreak{}|\allowbreak{}X\allowbreak{}\textless{}\allowbreak{}=\allowbreak{}c\allowbreak{}h\allowbreak{}i\allowbreak{}|\allowbreak{}|\allowbreak{}K\allowbreak{}h\allowbreak{}|\allowbreak{}|\allowbreak{}X\allowbreak{}.\allowbreak{}}
Direct algebra, with every term retained, yields

\begingroup\small\begin{equation*}
\adjustbox{max width=0.92\linewidth}{$\displaystyle \boxed{d_5(x)=3(1-\chi)
=\frac32(1+\sqrt{\mathscr D(x)})
+\sum_{i=1}^5(\tfrac32m_i-6t_i)x^i.}$}\tag{F35}
\end{equation*}\endgroup

Every last coefficient is nonnegative; 0\textless=chi\textless1/2 on
{[}0,alpha{]}. The correction w\_* increases on the open interval
because differentiation of
{(\allowbreak{}1\allowbreak{}-\allowbreak{}e\allowbreak{}l\allowbreak{}l\allowbreak{})\allowbreak{}w\allowbreak{}-\allowbreak{}(\allowbreak{}2\allowbreak{}/\allowbreak{}3\allowbreak{})\allowbreak{}w\allowbreak{}\textasciicircum{}\allowbreak{}2\allowbreak{}=\allowbreak{}d\allowbreak{}e\allowbreak{}l\allowbreak{}t\allowbreak{}a\allowbreak{}}
gives
{w\allowbreak{}'\allowbreak{}=\allowbreak{}(\allowbreak{}d\allowbreak{}e\allowbreak{}l\allowbreak{}t\allowbreak{}a\allowbreak{}'\allowbreak{}+\allowbreak{}e\allowbreak{}l\allowbreak{}l\allowbreak{}'\allowbreak{}w\allowbreak{})\allowbreak{}/\allowbreak{}s\allowbreak{}q\allowbreak{}r\allowbreak{}t\allowbreak{}(\allowbreak{}D\allowbreak{})\allowbreak{}\textgreater{}\allowbreak{}0\allowbreak{}.\allowbreak{}}
Therefore chi increases and d5 decreases to its positive endpoint.

Take an actual positive-energy physical eigenvector of H-E0. Division by
psi, then subtraction of its Haar scalar, gives a nonzero smooth h with

\begingroup\small\begin{equation*}
\adjustbox{max width=0.92\linewidth}{$\displaystyle (K-\lambda/\kappa)h=Q_H2\Gamma(v,h).$}\tag{F36}
\end{equation*}\endgroup

Smoothness implies sum\_j
{c\allowbreak{}\_\allowbreak{}j\allowbreak{}|\allowbreak{}|\allowbreak{}A\allowbreak{}\_\allowbreak{}j\allowbreak{}(\allowbreak{}h\allowbreak{})\allowbreak{}|\allowbreak{}|\allowbreak{}1\allowbreak{}\textless{}\allowbreak{}i\allowbreak{}n\allowbreak{}f\allowbreak{}i\allowbreak{}n\allowbreak{}i\allowbreak{}t\allowbreak{}y\allowbreak{}:\allowbreak{}}
Plancherel and Cauchy-\/-Schwarz with a sufficiently high original
Casimir power give this on the finite compact product. For
0\textless lambda/kappa\textless3,

\begingroup\small\begin{equation*}
\adjustbox{max width=\linewidth}{$\displaystyle \|(K-\lambda/\kappa)h\|_X
\ge(1-\lambda/(3\kappa))\|Kh\|_X.$}
\end{equation*}\endgroup

Combining with F35 proves lambda\textgreater=kappa d5. Values
lambda/kappa\textgreater=3 satisfy the same bound since d5\textless=3.
The complete compact physical spectral resolution consequently gives

\begingroup\small\begin{equation*}
\adjustbox{max width=0.92\linewidth}{$\displaystyle \boxed{\Delta_L\ge\kappa d_5(\xi),\qquad
q_{H-E_0}(\psi f)\ge\kappa d_5(\xi)\operatorname{Var}_{\psi^2}(f)}$}\tag{F37}
\end{equation*}\endgroup

for every L\textgreater=2, a\textgreater0 and
0\textless xi\textless=alpha, on the whole physical form domain. In
particular

\begingroup\small\begin{equation*}
\adjustbox{max width=0.92\linewidth}{$\displaystyle \begin{array}{c|c|c}
g^2\text{ at least}&\Delta_L/\kappa\text{ exceeds}&\|v-q_5\|_{\rm loc}\text{ below}\\\hline
37/10&1.6207&0.046581\\
15/4&1.6993&0.026352\\
4&1.9068&0.005752
\end{array}$}\tag{F38}
\end{equation*}\endgroup

The exact rational brackets, rather than these displayed rounded
consequences, are in
{p\allowbreak{}h\allowbreak{}y\allowbreak{}s\allowbreak{}i\allowbreak{}c\allowbreak{}a\allowbreak{}l\allowbreak{}\_\allowbreak{}r\allowbreak{}e\allowbreak{}t\allowbreak{}u\allowbreak{}r\allowbreak{}n\allowbreak{}.\allowbreak{}j\allowbreak{}s\allowbreak{}o\allowbreak{}n\allowbreak{}.\allowbreak{}}

The approximate exponential retains its complete operator defect:

\begingroup\small\begin{equation*}
\adjustbox{max width=0.92\linewidth}{$\displaystyle e^{-q_5}He^{q_5}
=\kappa[K-2\Gamma(q_5,\cdot)]
+\kappa[2M\xi-P_H\Gamma(q_5,q_5)]I-\kappa M_{KR_5}.$}\tag{F39}
\end{equation*}\endgroup

In particular, for the actual transformed operator Atilde,

\begingroup\small\begin{equation*}
\adjustbox{max width=0.92\linewidth}{$\displaystyle Q_H\widetilde Ah-\kappa KL_5h
=-2\kappa Q_H\Gamma(w,h),\qquad
\|\text{right side}\|_X\le\frac{2\kappa}{3}w_*\|Kh\|_X.$}\tag{F40}
\end{equation*}\endgroup

Thus the reference linearization and physical operator have their exact
comparison, including scalar and multiplication terms.

For source windows
{V\allowbreak{}\_\allowbreak{}N\allowbreak{}=\allowbreak{}s\allowbreak{}p\allowbreak{}a\allowbreak{}n\allowbreak{}\{\allowbreak{}R\allowbreak{}5\allowbreak{},\allowbreak{}J\allowbreak{}5\allowbreak{}R\allowbreak{}5\allowbreak{},\allowbreak{}.\allowbreak{}.\allowbreak{}.\allowbreak{},\allowbreak{}J\allowbreak{}5\allowbreak{}\textasciicircum{}\allowbreak{}N\allowbreak{}}
R5\}, the actual complexes are V\_N -\/-L5-\/-\textgreater{} V
-\/-0-\/-\textgreater0. Inclusion and identity give their transitions.
The map

\begingroup\small\begin{equation*}
\adjustbox{max width=0.92\linewidth}{$\displaystyle V_{N+1}/V_N\longrightarrow
\ker[V/L_5V_N\to V/L_5V_{N+1}],\quad [h]\mapsto[L_5h]$}\tag{F41}
\end{equation*}\endgroup

has inverse induced by L5\^{}\{-1\}: changing a representative by L5V\_N
changes its inverse by exactly V\_N. The finite primitive
sum\_(j=0)\^{}N J5\^{}jR5 has residual J5\^{}(N+1)R5. Its omitted norm
is at most
{e\allowbreak{}l\allowbreak{}l\allowbreak{}\textasciicircum{}\allowbreak{}(\allowbreak{}N\allowbreak{}+\allowbreak{}1\allowbreak{})\allowbreak{}d\allowbreak{}e\allowbreak{}l\allowbreak{}t\allowbreak{}a\allowbreak{}/\allowbreak{}(\allowbreak{}1\allowbreak{}-\allowbreak{}e\allowbreak{}l\allowbreak{}l\allowbreak{})\allowbreak{}.\allowbreak{}}
Both the original relation and its support label remain.

On centered physical functions let d\_X f=(X\_if)\_i with the original
range pairing kappa int rho sum
{c\allowbreak{}o\allowbreak{}n\allowbreak{}j\allowbreak{}u\allowbreak{}g\allowbreak{}a\allowbreak{}t\allowbreak{}e\allowbreak{}(\allowbreak{}v\allowbreak{}\_\allowbreak{}i\allowbreak{})\allowbreak{}w\allowbreak{}\_\allowbreak{}i\allowbreak{}.\allowbreak{}}
The inverse on its actual image is p\_X(d\_X f)=f. Its uniqueness
follows because all derivatives zero force a constant and centering
fixes that constant. The exact variational identity is

\begingroup\small\begin{equation*}
\adjustbox{max width=0.92\linewidth}{$\displaystyle \boxed{\|p_X\|^2=1/\Delta_L\le1/[\kappa d_5(\xi)].}$}\tag{F42}
\end{equation*}\endgroup

\subsection{F8. Sources and evidence}\label{f8.-sources-and-evidence}

The original LaTeX was located at repository commit
{a\allowbreak{}d\allowbreak{}d\allowbreak{}5\allowbreak{}9\allowbreak{}0\allowbreak{}f\allowbreak{}7\allowbreak{}8\allowbreak{}d\allowbreak{}0\allowbreak{}f\allowbreak{}6\allowbreak{}a\allowbreak{}1\allowbreak{}b\allowbreak{}d\allowbreak{}c\allowbreak{}c\allowbreak{}6\allowbreak{}e\allowbreak{}8\allowbreak{}3\allowbreak{}c\allowbreak{}c\allowbreak{}9\allowbreak{}e\allowbreak{}0\allowbreak{}3\allowbreak{}0\allowbreak{}1\allowbreak{}c\allowbreak{}f\allowbreak{}c\allowbreak{}9\allowbreak{}b\allowbreak{}0\allowbreak{}f\allowbreak{}5\allowbreak{}d\allowbreak{},\allowbreak{}}
{y\allowbreak{}a\allowbreak{}n\allowbreak{}g\allowbreak{}-\allowbreak{}m\allowbreak{}i\allowbreak{}l\allowbreak{}l\allowbreak{}s\allowbreak{}/\allowbreak{}c\allowbreak{}o\allowbreak{}n\allowbreak{}s\allowbreak{}o\allowbreak{}l\allowbreak{}i\allowbreak{}d\allowbreak{}a\allowbreak{}t\allowbreak{}i\allowbreak{}o\allowbreak{}n\allowbreak{}/\allowbreak{}2\allowbreak{}0\allowbreak{}2\allowbreak{}6\allowbreak{}0\allowbreak{}9\allowbreak{}1\allowbreak{}7\allowbreak{}/\allowbreak{}r\allowbreak{}e\allowbreak{}a\allowbreak{}d\allowbreak{}e\allowbreak{}r\allowbreak{}/\allowbreak{}y\allowbreak{}a\allowbreak{}n\allowbreak{}g\allowbreak{}\_\allowbreak{}m\allowbreak{}i\allowbreak{}l\allowbreak{}l\allowbreak{}s\allowbreak{}\_\allowbreak{}q\allowbreak{}u\allowbreak{}a\allowbreak{}r\allowbreak{}t\allowbreak{}i\allowbreak{}c\allowbreak{}\_\allowbreak{}c\allowbreak{}u\allowbreak{}b\allowbreak{}e\allowbreak{}\_\allowbreak{}r\allowbreak{}e\allowbreak{}a\allowbreak{}d\allowbreak{}e\allowbreak{}r\allowbreak{}.\allowbreak{}t\allowbreak{}e\allowbreak{}x\allowbreak{}.\allowbreak{}}
Its introductory original equations and six-appendix inventory were
compared with the mounted source chain. No exhaustive new review of that
whole LaTeX reader is claimed. The complete original fifth catalogue,
fourth input channel bounds and prior heat proof are retained unchanged
in this cumulative archive.

The explicit lower-order input polytopes occur in
{.\allowbreak{}.\allowbreak{}/\allowbreak{}2\allowbreak{}0\allowbreak{}2\allowbreak{}6\allowbreak{}0\allowbreak{}9\allowbreak{}1\allowbreak{}7\allowbreak{}-\allowbreak{}q\allowbreak{}u\allowbreak{}a\allowbreak{}r\allowbreak{}t\allowbreak{}i\allowbreak{}c\allowbreak{}-\allowbreak{}c\allowbreak{}u\allowbreak{}b\allowbreak{}e\allowbreak{}/\allowbreak{}c\allowbreak{}h\allowbreak{}a\allowbreak{}n\allowbreak{}n\allowbreak{}e\allowbreak{}l\allowbreak{}\_\allowbreak{}b\allowbreak{}o\allowbreak{}u\allowbreak{}n\allowbreak{}d\allowbreak{}s\allowbreak{}.\allowbreak{}p\allowbreak{}y\allowbreak{}}
and its
{Q\allowbreak{}U\allowbreak{}A\allowbreak{}R\allowbreak{}T\allowbreak{}I\allowbreak{}C\allowbreak{}\_\allowbreak{}S\allowbreak{}O\allowbreak{}U\allowbreak{}R\allowbreak{}C\allowbreak{}E\allowbreak{}.\allowbreak{}m\allowbreak{}d\allowbreak{}/\allowbreak{}P\allowbreak{}H\allowbreak{}Y\allowbreak{}S\allowbreak{}I\allowbreak{}C\allowbreak{}A\allowbreak{}L\allowbreak{}\_\allowbreak{}R\allowbreak{}E\allowbreak{}T\allowbreak{}U\allowbreak{}R\allowbreak{}N\allowbreak{}.\allowbreak{}m\allowbreak{}d\allowbreak{}}
proof. The full fifth source is
{.\allowbreak{}.\allowbreak{}/\allowbreak{}2\allowbreak{}0\allowbreak{}2\allowbreak{}6\allowbreak{}0\allowbreak{}9\allowbreak{}1\allowbreak{}7\allowbreak{}-\allowbreak{}f\allowbreak{}i\allowbreak{}f\allowbreak{}t\allowbreak{}h\allowbreak{}-\allowbreak{}s\allowbreak{}o\allowbreak{}u\allowbreak{}r\allowbreak{}c\allowbreak{}e\allowbreak{}/\allowbreak{}F\allowbreak{}I\allowbreak{}F\allowbreak{}T\allowbreak{}H\allowbreak{}\_\allowbreak{}S\allowbreak{}O\allowbreak{}U\allowbreak{}R\allowbreak{}C\allowbreak{}E\allowbreak{}.\allowbreak{}m\allowbreak{}d\allowbreak{}}
and its generated catalogue. This continuation\textquotesingle s new
algebra checks F11-\/-F25; the endpoint and physical-return arithmetic
are separately replayed. Analytic Banach-space and elliptic arguments
are written above, rather than represented as certified by finite matrix
tests.

Human antecedents: D. Schuette, Zheng Weihong and C. J. Hamer, The
Coupled Cluster Method in Hamiltonian Lattice Field Theory,
{a\allowbreak{}r\allowbreak{}X\allowbreak{}i\allowbreak{}v\allowbreak{}:\allowbreak{}h\allowbreak{}e\allowbreak{}p\allowbreak{}-\allowbreak{}l\allowbreak{}a\allowbreak{}t\allowbreak{}/\allowbreak{}9\allowbreak{}6\allowbreak{}0\allowbreak{}3\allowbreak{}0\allowbreak{}2\allowbreak{}6\allowbreak{}v\allowbreak{}1\allowbreak{},\allowbreak{}}
especially Sections 2-\/-5, provide the exponential vacuum, character
and Casimir framework. Their convention is connected by g\_C=2\^{}(3/4)g
and
{H\allowbreak{}\_\allowbreak{}o\allowbreak{}u\allowbreak{}r\allowbreak{}s\allowbreak{}=\allowbreak{}s\allowbreak{}q\allowbreak{}r\allowbreak{}t\allowbreak{}(\allowbreak{}2\allowbreak{})\allowbreak{}H\allowbreak{}\_\allowbreak{}C\allowbreak{}+\allowbreak{}2\allowbreak{}k\allowbreak{}a\allowbreak{}p\allowbreak{}p\allowbreak{}a\allowbreak{}}
xi M I, with their commutator derivative iX. P. Eymard,
L\textquotesingle algebre de Fourier d\textquotesingle un groupe
localement compact, Bull. Soc. Math. France 92 (1964),181-\/-236, is the
Fourier-algebra antecedent; the needed compact coefficient bounds are
proved here. No human reference is asserted to verify the new finite
catalogue or coupling threshold.

\section{Fifth-reference heat bounds and the complete complementary
physical
space}\label{fifth-reference-heat-bounds-and-the-complete-complementary-physical-space}

21 September 2026. This note applies the completed fifth-reference
correction in
{F\allowbreak{}I\allowbreak{}F\allowbreak{}T\allowbreak{}H\allowbreak{}\_\allowbreak{}R\allowbreak{}E\allowbreak{}F\allowbreak{}E\allowbreak{}R\allowbreak{}E\allowbreak{}N\allowbreak{}C\allowbreak{}E\allowbreak{}.\allowbreak{}m\allowbreak{}d\allowbreak{}}
to the original Yang--Mills heat, plaquette moments and complementary
minimization. The improvement comes from evaluating that source, not
from changing a physical norm. It uses the unchanged exact heat
coefficients through degree four and constructs a new box-independent
analytic remainder. All exterior-volume, observation-rank and
physical-time scopes are stated explicitly.

\subsection{H1. Complete coefficient generator and original
mean}\label{h1.-complete-coefficient-generator-and-original-mean}

Retain F1-\/-F42 of the companion, in particular kappa=2g\^{}2/a,
xi=1/(4g\^{}4), rho=psi\^{}2 and A=H-E0. Fix a complex radius
r\textless alpha. The fifth construction gives a local analytic source
v(zeta) and a drift D\_v=2Gamma(v,.), satisfying

\begingroup\small\begin{equation*}
\adjustbox{max width=0.92\linewidth}{$\displaystyle \|D_vh\|_X\le\chi(r)\|Kh\|_X,\quad
\chi(r)=1-d_5(r)/3<1/2.$}\tag{H1}
\end{equation*}\endgroup

Let X0 be the nonconstant physical Fourier space with norm sum
\textbar\textbar A\_j\textbar\textbar1, and Y0 its domain sum
{c\allowbreak{}\_\allowbreak{}j\allowbreak{}|\allowbreak{}|\allowbreak{}A\allowbreak{}\_\allowbreak{}j\allowbreak{}|\allowbreak{}|\allowbreak{}1\allowbreak{}\textless{}\allowbreak{}i\allowbreak{}n\allowbreak{}f\allowbreak{}i\allowbreak{}n\allowbreak{}i\allowbreak{}t\allowbreak{}y\allowbreak{}.\allowbreak{}}
Define the actual maps

\begingroup\small\begin{equation*}
\adjustbox{max width=0.92\linewidth}{$\displaystyle T=Q_HD_vK^{-1}:X_0\to X_0,\quad p=P_HD_vK^{-1}:X_0\to\mathbb C,
\quad S=(I-T)^{-1},\quad
\mu(F)=P_HF+pSQ_HF.$}\tag{H2}
\end{equation*}\endgroup

The complete scalar-and-nonconstant output obeys
{|\allowbreak{}p\allowbreak{}y\allowbreak{}|\allowbreak{}+\allowbreak{}|\allowbreak{}|\allowbreak{}T\allowbreak{}y\allowbreak{}|\allowbreak{}|\allowbreak{}X\allowbreak{}\textless{}\allowbreak{}=\allowbreak{}c\allowbreak{}h\allowbreak{}i\allowbreak{}|\allowbreak{}|\allowbreak{}y\allowbreak{}|\allowbreak{}|\allowbreak{}X\allowbreak{}.\allowbreak{}}
Thus \textbar\textbar S\textbar\textbar\textless=1/(1-chi),
\textbar\textbar mu\textbar\textbar\textless=1, and
{|\allowbreak{}(\allowbreak{}m\allowbreak{}u\allowbreak{}-\allowbreak{}P\allowbreak{}\_\allowbreak{}H\allowbreak{})\allowbreak{}Q\allowbreak{}\_\allowbreak{}H\allowbreak{}F\allowbreak{}|\allowbreak{}\textless{}\allowbreak{}=\allowbreak{}c\allowbreak{}h\allowbreak{}i\allowbreak{}|\allowbreak{}|\allowbreak{}Q\allowbreak{}\_\allowbreak{}H\allowbreak{}F\allowbreak{}|\allowbreak{}|\allowbreak{}X\allowbreak{}/\allowbreak{}(\allowbreak{}1\allowbreak{}-\allowbreak{}c\allowbreak{}h\allowbreak{}i\allowbreak{})\allowbreak{}.\allowbreak{}}
The last bound and chi\textless=1/2 prove the full mean norm, using the
additive original
{c\allowbreak{}o\allowbreak{}n\allowbreak{}s\allowbreak{}t\allowbreak{}a\allowbreak{}n\allowbreak{}t\allowbreak{}/\allowbreak{}n\allowbreak{}o\allowbreak{}n\allowbreak{}c\allowbreak{}o\allowbreak{}n\allowbreak{}s\allowbreak{}t\allowbreak{}a\allowbreak{}n\allowbreak{}t\allowbreak{}}
coefficient norm. The exact Poisson equation is

\begingroup\small\begin{equation*}
\adjustbox{max width=0.92\linewidth}{$\displaystyle (K-D_v)K^{-1}SQ_HF=F-\mu(F).$}\tag{H3}
\end{equation*}\endgroup

This constructs the scalar as well as the primitive. Integration by
parts against exp(2v) gives int exp(2v)F=mu(F)int exp(2v). The latter
integral cannot vanish: a vanishing value would make the left side zero
for every physical Fourier polynomial, whose uniform closure contains
the conjugate of the nonzero smooth gauge-invariant exp(2v). Their
paired integral would then be zero. Consequently

\begingroup\small\begin{equation*}
\adjustbox{max width=0.92\linewidth}{$\displaystyle \mu(F)=\frac{\int e^{2v}F\,dU}{\int e^{2v}\,dU}.$}\tag{H4}
\end{equation*}\endgroup

On real coupling this is the original vacuum expectation. The complex
formula keeps its full nonzero scalar denominator.

The coefficient operator L=(I-T)K on Y0 is closed, since I-T has a
bounded inverse and K is closed. Finite Fourier polynomials give a dense
domain. For h\textgreater0 the factorization of I+hL through I+hK has
remaining factor
{I\allowbreak{}-\allowbreak{}h\allowbreak{}Q\allowbreak{}\_\allowbreak{}H\allowbreak{}D\allowbreak{}\_\allowbreak{}v\allowbreak{}(\allowbreak{}I\allowbreak{}+\allowbreak{}h\allowbreak{}K\allowbreak{})\allowbreak{}\textasciicircum{}\allowbreak{}\{\allowbreak{}-\allowbreak{}1\allowbreak{}\}\allowbreak{},\allowbreak{}}
whose norm difference from I is at most chi. It is therefore onto with
bounded inverse. For f in Y0,

\begingroup\small\begin{equation*}
\adjustbox{max width=0.92\linewidth}{$\displaystyle \|(I+hL)f\|_X\ge\|f\|_X+h(1-\chi)\|Kf\|_X
\ge[1+3h(1-\chi)]\|f\|_X.$}\tag{H5}
\end{equation*}\endgroup

Here the identity
\textbar\textbar(I+hK)f\textbar\textbar=\textbar\textbar f\textbar\textbar+h\textbar\textbar Kf\textbar\textbar{}
follows coefficient by coefficient. The same derivative estimate proves
dissipativity after shifting by 3(1-chi); range surjectivity at a
sufficiently large positive parameter follows from the displayed
factorization. The
{H\allowbreak{}i\allowbreak{}l\allowbreak{}l\allowbreak{}e\allowbreak{}-\allowbreak{}-\allowbreak{}Y\allowbreak{}o\allowbreak{}s\allowbreak{}i\allowbreak{}d\allowbreak{}a\allowbreak{}/\allowbreak{}L\allowbreak{}u\allowbreak{}m\allowbreak{}e\allowbreak{}r\allowbreak{}-\allowbreak{}-\allowbreak{}P\allowbreak{}h\allowbreak{}i\allowbreak{}l\allowbreak{}l\allowbreak{}i\allowbreak{}p\allowbreak{}s\allowbreak{}}
generation theorem thus gives

\begingroup\small\begin{equation*}
\adjustbox{max width=0.92\linewidth}{$\displaystyle E(\tau)=e^{-\tau L},\qquad \|E(\tau)\|\le e^{-3(1-\chi)\tau}.$}\tag{H6}
\end{equation*}\endgroup

Both the domain and actual range were verified, rather than inferring a
semigroup bound from a single inverse. Analytic parameter dependence
follows from the locally analytic resolvents and their exponential Euler
powers: each power is a positive scalar gamma-time average of E, has the
same contraction bound, and converges strongly by its mean and variance.
Bounded weak holomorphic convergence, tested against every continuous
linear functional, gives the analytic continuation on every smaller
source disk.

The original full heat is

\begingroup\small\begin{equation*}
\adjustbox{max width=0.92\linewidth}{$\displaystyle \mathcal T(\tau)F=E(\tau)Q_HF+\mu(F)-\mu(E(\tau)Q_HF).$}\tag{H7}
\end{equation*}\endgroup

It solves the original transformed equation K-D\_v, preserves the actual
mean, and on real coupling agrees with the physical Hilbert semigroup by
uniqueness. The physical time map and inverse are tau=kappa t and
t=tau/kappa.

\subsection{H2. New box-independent heat
circle}\label{h2.-new-box-independent-heat-circle}

For original marked faces define

\begingroup\small\begin{equation*}
\adjustbox{max width=0.92\linewidth}{$\displaystyle r_p=(W_p-\mu W_p)\psi,\quad
\widehat C_{pq}(\tau;\xi)=\langle r_p,e^{-\tau A/\kappa}r_q\rangle.$}\tag{H8}
\end{equation*}\endgroup

At complex parameter use its bilinear analytic continuation
mu{[}(T(tau)Wp-mu Wp)Wq{]}. For arbitrary original face coefficients
\textbar z\_q\textbar\textless=1 put G\_z=sum\_q z\_qWq. At most four
faces meet an edge. Their fundamental coefficient norm 8 and edge spin
1/2 give

\begingroup\small\begin{equation*}
\adjustbox{max width=0.92\linewidth}{$\displaystyle \|\Gamma(h,G_z)\|_X\le32\|Kh\|_X.$}\tag{H9}
\end{equation*}\endgroup

For the original Haar scalar there is a sharper exact selection rule. A
physical representation block with spin 1/2 on one elementary face and
zero elsewhere is the one-dimensional line C Wp. Its Casimir is 3,
coefficient norm is 8\textbar a\_p\textbar, and Haar norm of Wp is 1.
Integration by parts and Haar orthogonality therefore give

\begingroup\small\begin{equation*}
\adjustbox{max width=0.92\linewidth}{$\displaystyle |P_H\Gamma(h,G_z)|\le\frac18\|Kh\|_X.$}\tag{H10}
\end{equation*}\endgroup

Take h\_tau=K\^{}\{-1\}S E(tau)Wp. Equation H3, followed by integration
by parts in H4, proves

\begingroup\small\begin{equation*}
\adjustbox{max width=\linewidth}{$\displaystyle \sum_qz_q\widehat C_{pq}(\tau)=\mu\Gamma(h_\tau,G_z),\quad
\|Kh_\tau\|_X\le\frac8{1-\chi}e^{-3(1-\chi)\tau}.$}
\end{equation*}\endgroup

Retaining the Haar contribution (H10) and the full returned-mean
contribution in H2 proves

\begingroup\small\begin{equation*}
\adjustbox{max width=0.92\linewidth}{$\displaystyle \boxed{\|\widehat C(\tau;\zeta)\|_{\rm row}
\le\frac{1+255\chi}{(1-\chi)^2}e^{-3(1-\chi)\tau}.}$}\tag{H11}
\end{equation*}\endgroup

The row norm is max\_p sum\_q\textbar C\_pq\textbar, obtained exactly by
maximizing over the phases z\_q. Transpose symmetry extends from real
coupling by analytic continuation, so columns have the same bound. No
face count M enters this estimate.

The exact source calculation F28-\/-F35 verifies

\begingroup\small\begin{equation*}
\adjustbox{max width=0.92\linewidth}{$\displaystyle \boxed{R=1/55<\alpha,\quad d_5(R)>13/8,\quad
\chi(R)<11/24,\quad C=\frac{1+255(11/24)}{(1-11/24)^2}
=\frac{67896}{169}.}$}\tag{H12}
\end{equation*}\endgroup

This circle strictly exceeds the saved fourth-reference endpoint
0.018104972231644127076, and the preceding volume-heat radius 3/256. The
inequalities are certified by integer-square and rational polynomial
bounds in
{p\allowbreak{}h\allowbreak{}y\allowbreak{}s\allowbreak{}i\allowbreak{}c\allowbreak{}a\allowbreak{}l\allowbreak{}\_\allowbreak{}r\allowbreak{}e\allowbreak{}t\allowbreak{}u\allowbreak{}r\allowbreak{}n\allowbreak{}.\allowbreak{}j\allowbreak{}s\allowbreak{}o\allowbreak{}n\allowbreak{}.\allowbreak{}}

The original link-center map U\_(n,i) -\textgreater{}
{(\allowbreak{}-\allowbreak{}1\allowbreak{})\allowbreak{}\textasciicircum{}\allowbreak{}(\allowbreak{}s\allowbreak{}u\allowbreak{}m\allowbreak{}\_\allowbreak{}(\allowbreak{}j\allowbreak{}\textless{}\allowbreak{}i\allowbreak{})\allowbreak{}n\allowbreak{}\_\allowbreak{}j\allowbreak{})\allowbreak{}}
U\_(n,i) preserves Haar, K and the full gauge action and sends every Wp
to -Wp. Therefore the homogeneous connected two-mark heat is even.
Applying Cauchy\textquotesingle s formula to (H11) on the same circle
gives

\begingroup\small\begin{equation*}
\adjustbox{max width=0.92\linewidth}{$\displaystyle \boxed{
\|\widehat C(\tau;\xi)-C_0(\tau)-\xi^2C_2(\tau)-\xi^4C_4(\tau)\|_{\rm row}
\le\frac{67896}{169}e^{-13\tau/8}
\frac{(55|\xi|)^6}{1-(55|\xi|)^2},\quad |\xi|<1/55.}$}\tag{H13}
\end{equation*}\endgroup

This holds for every original L\textgreater=2 and every
tau\textgreater=0. C0,C2,C4 are the complete unchanged inherited heat
coefficient matrices; the original source/heat recurrence fixes them
uniquely. The new result evaluates an infinite tail beyond those
coefficients; it does not relabel a new finite heat coefficient as
calculated.

\subsection{H3. Original moments and a signed leakage
remainder}\label{h3.-original-moments-and-a-signed-leakage-remainder}

Let R\_col have the original columns r\_p and set Phi=kappa
A\^{}\{-1\}R\_col. Keep

\begingroup\small\begin{equation*}
\adjustbox{max width=0.92\linewidth}{$\displaystyle G_0=R_{\rm col}^*R_{\rm col},\quad
G_1=\kappa R_{\rm col}^*A^{-1}R_{\rm col},\quad
G_2=\Phi^*\Phi,\quad
E=\Phi^*A\Phi=\kappa G_1,\quad
K_o=\kappa^{-1}R_{\rm col}^*AR_{\rm col}.$}\tag{H14}
\end{equation*}\endgroup

The spectral theorem gives
{G\allowbreak{}\_\allowbreak{}k\allowbreak{}=\allowbreak{}i\allowbreak{}n\allowbreak{}t\allowbreak{}\_\allowbreak{}0\allowbreak{}\textasciicircum{}\allowbreak{}i\allowbreak{}n\allowbreak{}f\allowbreak{}i\allowbreak{}n\allowbreak{}i\allowbreak{}t\allowbreak{}y\allowbreak{}}
tau\^{}(k-1)Chat(tau)d tau/(k-1)! for k\textgreater=1. Thus (H13) has
moment prefactor C/d\^{}k, d=13/8, retaining every power of kappa in
H14. G0 uses tau=0.

The complete kinetic function is
{\textless{}\allowbreak{}G\allowbreak{}a\allowbreak{}m\allowbreak{}m\allowbreak{}a\allowbreak{}(\allowbreak{}W\allowbreak{}p\allowbreak{},\allowbreak{}W\allowbreak{}q\allowbreak{})\allowbreak{}\textgreater{}\allowbreak{}r\allowbreak{}h\allowbreak{}o\allowbreak{}.\allowbreak{}}
A diagonal has coefficient norm at most30, from
{G\allowbreak{}a\allowbreak{}m\allowbreak{}m\allowbreak{}a\allowbreak{}(\allowbreak{}W\allowbreak{}p\allowbreak{},\allowbreak{}W\allowbreak{}p\allowbreak{})\allowbreak{}=\allowbreak{}4\allowbreak{}-\allowbreak{}W\allowbreak{}p\allowbreak{}\textasciicircum{}\allowbreak{}2\allowbreak{}=\allowbreak{}3\allowbreak{}-\allowbreak{}c\allowbreak{}h\allowbreak{}i\allowbreak{}\_\allowbreak{}1\allowbreak{}}
and \textbar\textbar chi\_1\textbar\textbar X=27. An adjacent
off-diagonal has norm at most48. Every face has at most twelve adjacent
faces; every other entry is zero. Since
\textbar\textbar mu\textbar\textbar\textless=1, its full complex row
bound is 606. It consequently has the same even tail factor, with
prefactor606.

The actual nonnegative leakage trial is

\begingroup\small\begin{equation*}
\adjustbox{max width=0.92\linewidth}{$\displaystyle J=(R_{\rm col}-3\Phi)^*(R_{\rm col}-3\Phi)=G_0-6G_1+9G_2.$}\tag{H15}
\end{equation*}\endgroup

Keep the sign of the time weight before estimating its remainder:

\begingroup\small\begin{equation*}
\adjustbox{max width=0.92\linewidth}{$\displaystyle J=\widehat C(0)+\int_0^\infty(9\tau-6)\widehat C(\tau)\,d\tau.$}\tag{H16}
\end{equation*}\endgroup

Splitting the scalar integral at tau=2/3 gives the exact expression

\begingroup\small\begin{equation*}
\adjustbox{max width=0.92\linewidth}{$\displaystyle 1+\int_0^\infty|9\tau-6|e^{-d\tau}d\tau
=1+\frac6d-\frac9{d^2}+\frac{18}{d^2}e^{-2d/3}.$}\tag{H17}
\end{equation*}\endgroup

For d=13/8, the even 120th Taylor sum of exp(-13/12) is less
than339/1000. Taylor\textquotesingle s integral remainder is negative at
that order, giving a rigorous upper bound for the exponential. Hence the
complete leakage-tail prefactor is at most

\begingroup\small\begin{equation*}
\adjustbox{max width=0.92\linewidth}{$\displaystyle \boxed{C_J=C\left[1+\frac6d-\frac9{d^2}+\frac{18(339/1000)}{d^2}\right]
=\frac{5156090136}{3570125}.}$}\tag{H18}
\end{equation*}\endgroup

This is smaller than C(1+6/d+9/d\^{}2). All signed time weights and
mixed moment entries remain in H15-\/-H18.

The complete inherited original-support coefficient bounds are

{\def\LTcaptype{none} % do not increment counter
\begin{longtable}[]{@{}lrr@{}}
\toprule\noalign{}
Matrix & degree two row bound & degree four row bound \\
\midrule\noalign{}
\endhead
\bottomrule\noalign{}
\endlastfoot
G0 & 149/468 &
{2\allowbreak{}3\allowbreak{}6\allowbreak{}1\allowbreak{}9\allowbreak{}9\allowbreak{}4\allowbreak{}0\allowbreak{}7\allowbreak{}3\allowbreak{}/\allowbreak{}4\allowbreak{}6\allowbreak{}9\allowbreak{}1\allowbreak{}4\allowbreak{}9\allowbreak{}4\allowbreak{}0\allowbreak{}8\allowbreak{}0\allowbreak{}} \\
G1 & 97/468 &
{3\allowbreak{}7\allowbreak{}6\allowbreak{}8\allowbreak{}8\allowbreak{}2\allowbreak{}6\allowbreak{}9\allowbreak{}1\allowbreak{}/\allowbreak{}9\allowbreak{}3\allowbreak{}8\allowbreak{}2\allowbreak{}9\allowbreak{}8\allowbreak{}8\allowbreak{}1\allowbreak{}6\allowbreak{}} \\
G2 & 73313/657072 &
{9\allowbreak{}8\allowbreak{}2\allowbreak{}0\allowbreak{}6\allowbreak{}9\allowbreak{}7\allowbreak{}1\allowbreak{}8\allowbreak{}9\allowbreak{}6\allowbreak{}3\allowbreak{}8\allowbreak{}3\allowbreak{}3\allowbreak{}/\allowbreak{}3\allowbreak{}9\allowbreak{}1\allowbreak{}9\allowbreak{}1\allowbreak{}8\allowbreak{}0\allowbreak{}3\allowbreak{}2\allowbreak{}4\allowbreak{}5\allowbreak{}5\allowbreak{}0\allowbreak{}4\allowbreak{}0\allowbreak{}0\allowbreak{}} \\
Ko & 187/468 &
{5\allowbreak{}8\allowbreak{}6\allowbreak{}6\allowbreak{}6\allowbreak{}8\allowbreak{}4\allowbreak{}2\allowbreak{}1\allowbreak{}/\allowbreak{}1\allowbreak{}5\allowbreak{}6\allowbreak{}3\allowbreak{}8\allowbreak{}3\allowbreak{}1\allowbreak{}3\allowbreak{}6\allowbreak{}0\allowbreak{}} \\
J & 6779/73008 &
{1\allowbreak{}5\allowbreak{}2\allowbreak{}3\allowbreak{}3\allowbreak{}8\allowbreak{}6\allowbreak{}7\allowbreak{}4\allowbreak{}0\allowbreak{}0\allowbreak{}5\allowbreak{}9\allowbreak{}8\allowbreak{}9\allowbreak{}/\allowbreak{}4\allowbreak{}3\allowbreak{}5\allowbreak{}4\allowbreak{}6\allowbreak{}4\allowbreak{}4\allowbreak{}8\allowbreak{}0\allowbreak{}5\allowbreak{}0\allowbreak{}5\allowbreak{}6\allowbreak{}0\allowbreak{}0\allowbreak{}} \\
\end{longtable}
}

Their actual absolute-motif assembly retains all 199 original supports
and559 marked contributions in the preceding volume directory. Boundary
rows use subsets; no claim that a signed bulk row is the worst boundary
row is used.

For x=xi and
{u\allowbreak{}=\allowbreak{}(\allowbreak{}5\allowbreak{}5\allowbreak{}x\allowbreak{})\allowbreak{}\textasciicircum{}\allowbreak{}6\allowbreak{}/\allowbreak{}[\allowbreak{}1\allowbreak{}-\allowbreak{}(\allowbreak{}5\allowbreak{}5\allowbreak{}x\allowbreak{})\allowbreak{}\textasciicircum{}\allowbreak{}2\allowbreak{}]\allowbreak{},\allowbreak{}}
define the explicit widths

\begingroup\small\begin{equation*}
\adjustbox{max width=0.92\linewidth}{$\displaystyle \begin{aligned}
w_k&=B_{k,2}x^2+B_{k,4}x^4+(C/d^k)u\quad(k=0,1,2),\\
w_K&=B_{K,2}x^2+B_{K,4}x^4+606u,\\
\beta&=B_{J,2}x^2+B_{J,4}x^4+C_Ju.
\end{aligned}$}\tag{H19}
\end{equation*}\endgroup

The original real symmetric matrices therefore satisfy

\begingroup\small\begin{equation*}
\adjustbox{max width=0.92\linewidth}{$\displaystyle \|G_k-3^{-k}I\|\le w_k,\quad\|K_o-3I\|\le w_K,
\qquad0\preceq J\preceq\beta I.$}\tag{H20}
\end{equation*}\endgroup

The operator estimate follows from both row and column bounds. Each
width increases with x. Their exact rational values at the original
coupling benchmarks, with all square-root enclosures for d5, are in
{g\allowbreak{}e\allowbreak{}n\allowbreak{}e\allowbreak{}r\allowbreak{}a\allowbreak{}t\allowbreak{}e\allowbreak{}d\allowbreak{}/\allowbreak{}p\allowbreak{}h\allowbreak{}y\allowbreak{}s\allowbreak{}i\allowbreak{}c\allowbreak{}a\allowbreak{}l\allowbreak{}\_\allowbreak{}r\allowbreak{}e\allowbreak{}t\allowbreak{}u\allowbreak{}r\allowbreak{}n\allowbreak{}.\allowbreak{}j\allowbreak{}s\allowbreak{}o\allowbreak{}n\allowbreak{}.\allowbreak{}}

\subsection{H4. Whole physical observation and complementary
coupling}\label{h4.-whole-physical-observation-and-complementary-coupling}

Put l\_k=3\^{}(-k)-w\_k, u\_k=3\^{}(-k)+w\_k for k=0,1,2, and
u\_K=3+w\_K. Each benchmark below has l\_k\textgreater0. Thus R\_col and
Phi are injective and have closed ranges, with the original Grams H14.
The actual state projection is

\begingroup\small\begin{equation*}
\adjustbox{max width=\linewidth}{$\displaystyle P_R=R_{\rm col}G_0^{-1}R_{\rm col}^*,\quad
\|P_R\Phi c\|^2=c^*G_1G_0^{-1}G_1c.$}
\end{equation*}\endgroup

Consequently

\begingroup\small\begin{equation*}
\adjustbox{max width=0.92\linewidth}{$\displaystyle \frac{\|P_R\Phi c\|^2}{\|\Phi c\|^2}
\ge\frac{l_1^2}{u_0u_2}\quad(c\ne0).$}\tag{H21}
\end{equation*}\endgroup

This is observation of the actual inverse-energy vectors, with both
original state metrics retained.

For the complete physical complement define

\begingroup\small\begin{equation*}
\adjustbox{max width=0.92\linewidth}{$\displaystyle P=\Phi G_2^{-1}\Phi^*,\quad Q=I-P,\quad B_\Phi=QR_{\rm col}.$}\tag{H22}
\end{equation*}\endgroup

Completing the original state-norm square gives

\begingroup\small\begin{equation*}
\adjustbox{max width=0.92\linewidth}{$\displaystyle B_\Phi^*B_\Phi=G_0-G_1G_2^{-1}G_1\preceq J\preceq\beta I.$}\tag{H23}
\end{equation*}\endgroup

The comparison with J follows by inserting the specific admissible
coefficient3c into the same state minimum. Every centered physical
form-domain vector has the decomposition Phi c+h, h in QH0. Its complete
pairing and energy are

\begingroup\small\begin{equation*}
\adjustbox{max width=0.92\linewidth}{$\displaystyle \begin{aligned}
\|\Phi c+h\|^2&=c^*G_2c+\|h\|^2,\\
q_A(\Phi c+h)&=\kappa c^*G_1c
+2\kappa\operatorname{Re}\langle B_\Phi c,h\rangle+q_A(h).
\end{aligned}$}\tag{H24}
\end{equation*}\endgroup

Take d\_ph as any displayed rational lower endpoint for d5(x); then
A\textgreater=kappa d\_ph on H0 by F37. The product A Phi=kappa R\_col
is bounded. Hence AP is bounded, PA has its bounded adjoint extension,
and QAP+PAQ is bounded self-adjoint. Subtraction of this operator from A
leaves a self-adjoint block-diagonal operator on Dom(A). Its Q
restriction is

\begingroup\small\begin{equation*}
\adjustbox{max width=0.92\linewidth}{$\displaystyle D_Q=QAQ,\quad\operatorname{Dom}D_Q=\operatorname{Dom}A\cap Q\mathcal H_0,
\quad D_Q\succeq\kappa d_{\rm ph}I.$}\tag{H25}
\end{equation*}\endgroup

This proves its operator domain and inverse on the complete
complementary space, rather than just a trial subspace. Direct
substitution minimizes H24 at

\begingroup\small\begin{equation*}
\adjustbox{max width=\linewidth}{$\displaystyle h_{\min}=-\kappa D_Q^{-1}B_\Phi c.$}
\end{equation*}\endgroup

The restored columns and metrics are exactly

\begingroup\small\begin{equation*}
\adjustbox{max width=0.92\linewidth}{$\displaystyle \begin{aligned}
\Phi_{\rm full}&=\Phi-\kappa D_Q^{-1}B_\Phi,\\
E_{\rm full}&=\kappa G_1-\kappa^2B_\Phi^*D_Q^{-1}B_\Phi,\\
G_{\rm full}&=G_2+\kappa^2B_\Phi^*D_Q^{-2}B_\Phi.
\end{aligned}$}\tag{H26}
\end{equation*}\endgroup

Define the explicit positive numbers

\begingroup\small\begin{equation*}
\adjustbox{max width=0.92\linewidth}{$\displaystyle \eta_E=\frac{\beta}{d_{\rm ph}l_1},\quad
\eta_G=\frac{\beta}{d_{\rm ph}^2l_2}.$}\tag{H27}
\end{equation*}\endgroup

Then

\begingroup\small\begin{equation*}
\adjustbox{max width=0.92\linewidth}{$\displaystyle (1-\eta_E)E\preceq E_{\rm full}\preceq E,\qquad
G_2\preceq G_{\rm full}\preceq(1+\eta_G)G_2.$}\tag{H28}
\end{equation*}\endgroup

The full mixed form H24 also lies between (1-sqrt(eta\_E)) and
(1+sqrt(eta\_E)) times kappa c*G1c+q\_A(h), by Cauchy-\/-Schwarz in the
actual D\_Q energy. The cross term has been bounded, not discarded.

The state minimum section P\_R Phi differs from the energy minimum
section Phi by h\_R=(I-P\_R)Phi. The original equation A Phi=kappa
R\_col gives q\_A(Phi c,h\_Rd)=0. Expanding both mixed terms gives

\begingroup\small\begin{equation*}
\adjustbox{max width=0.92\linewidth}{$\displaystyle h_R^*Ah_R=\kappa[G_1G_0^{-1}K_oG_0^{-1}G_1-G_1]
\preceq\left(\frac{u_1u_K}{l_0^2}-1\right)E.$}\tag{H29}
\end{equation*}\endgroup

\subsection{H5. Explicit improved
domains}\label{h5.-explicit-improved-domains}

The rational evaluation of H19-\/-H29 proves, simultaneously for all
original boxes L\textgreater=2 and a\textgreater0:

{\def\LTcaptype{none} % do not increment counter
\begin{longtable}[]{@{}
  >{\raggedright\arraybackslash}p{(\linewidth - 6\tabcolsep) * \real{0.2000}}
  >{\raggedleft\arraybackslash}p{(\linewidth - 6\tabcolsep) * \real{0.2667}}
  >{\raggedleft\arraybackslash}p{(\linewidth - 6\tabcolsep) * \real{0.2667}}
  >{\raggedleft\arraybackslash}p{(\linewidth - 6\tabcolsep) * \real{0.2667}}@{}}
\toprule\noalign{}
\begin{minipage}[b]{\linewidth}\raggedright
Original coupling
\end{minipage} & \begin{minipage}[b]{\linewidth}\raggedleft
Observed state fraction exceeds
\end{minipage} & \begin{minipage}[b]{\linewidth}\raggedleft
Complementary energy loss below
\end{minipage} & \begin{minipage}[b]{\linewidth}\raggedleft
Restored state increase below
\end{minipage} \\
\midrule\noalign{}
\endhead
\bottomrule\noalign{}
\endlastfoot
g²\textgreater=10 & 977/1000 & 11/1000 & 12/1000 \\
g²\textgreater=12 & 997/1000 & 12/10000 & 12/10000 \\
g²\textgreater=25/2 & 998/1000 & 1/1000 & 1/1000 \\
g²\textgreater=13 & 999/1000 & 1/2000 & 1/2000 \\
g²\textgreater=16 & 9999/10000 & 1/25000 & 1/25000 \\
\end{longtable}
}

``Loss below eta'' in this table means E\_full\textgreater(1-eta)E;
``increase below eta'' means G\_full\textless(1+eta)G2. All identities
and quotients are the original H14,H26, not substituted metrics. At
g²=13, the exact diagnostic values of the three respective bounds are
approximately0.99904459483,0.00043585425,0.00045023706; the conclusions
in the table use strict rational tests. For g²\textgreater=16 the
corresponding original section-energy change H29 is below1/20000.

For example at g²\textgreater=13 the full original Gram bounds include

\begingroup\small\begin{equation*}
\adjustbox{max width=0.92\linewidth}{$\displaystyle \|G_0-I\|<119\cdot10^{-6},\quad
\|G_1-I/3\|<73\cdot10^{-6},\quad
\|G_2-I/9\|<45\cdot10^{-6},\quad
\|K_o-3I\|<178\cdot10^{-6}.$}\tag{H30}
\end{equation*}\endgroup

The physical gap used there has d\_ph=29047/10000. Every estimate for
larger g² follows because xi decreases, the widths increase with xi, and
d5 decreases with xi.

\subsection{H6. Support quotients, heat horizon, and spatial
limit}\label{h6.-support-quotients-heat-horizon-and-spatial-limit}

For original face sets F subset G, keep V\_F=Phi(C\^{}F), U\_F=A V\_F.
The cochain windows are V\_F -\/-A-\/-\textgreater{} H0
-\/-0-\/-\textgreater0. The exact transported-kernel isomorphism is

\begingroup\small\begin{equation*}
\adjustbox{max width=0.92\linewidth}{$\displaystyle V_G/V_F\longrightarrow\ker[\mathcal H_0/U_F\to\mathcal H_0/U_G],
\qquad[h]\mapsto[Ah].$}\tag{H31}
\end{equation*}\endgroup

Its inverse is induced by the actual A\^{}\{-1\}; changing a
representative by U\_F changes its inverse by V\_F. For J=G minus F, its
original energy quotient is the attained minimum

\begingroup\small\begin{equation*}
\adjustbox{max width=0.92\linewidth}{$\displaystyle Q_E(F,G)=E_{JJ}-E_{JF}E_{FF}^{-1}E_{FJ},\qquad
c_F^{\min}=-E_{FF}^{-1}E_{FJ}c_J.$}\tag{H32}
\end{equation*}\endgroup

Every mixed block and the removed primitive remain in the formula. The
finite-time columns

\begingroup\small\begin{equation*}
\adjustbox{max width=\linewidth}{$\displaystyle \Phi_T=(I-e^{-TA})\Phi=\kappa\int_0^Te^{-tA}R_{\rm col}\,dt$}
\end{equation*}\endgroup

satisfy A Phi\_T=kappa R\_col-kappa e\^{}\{-TA\}R\_col. With
epsilon=e\^{}(-kappa d\_ph T), the original state and energy Grams lie
between (1-epsilon)\^{}2 and1 times the original Grams. This follows
spectrally from A\textgreater=kappa d\_ph and commutation with exp(-TA).
Taking infima over the identical coefficient fibers in H32 preserves
these inequalities.

For a prescribed finite signed determinant return sum\_j a\_j logdet
Q\_j of total absolute rank B=sum\_j\textbar a\_j\textbar rank Q\_j, the
error is at most2B{[}-log(1-epsilon){]}. A finite physical choice is

\begingroup\small\begin{equation*}
\adjustbox{max width=0.92\linewidth}{$\displaystyle \boxed{T=(\kappa d_{\rm ph})^{-1}\log(1+2B/\eta),}$}\tag{H33}
\end{equation*}\endgroup

which gives error\textless=eta by
{-\allowbreak{}l\allowbreak{}o\allowbreak{}g\allowbreak{}(\allowbreak{}1\allowbreak{}-\allowbreak{}e\allowbreak{}p\allowbreak{}s\allowbreak{}i\allowbreak{}l\allowbreak{}o\allowbreak{}n\allowbreak{})\allowbreak{}\textless{}\allowbreak{}=\allowbreak{}e\allowbreak{}p\allowbreak{}s\allowbreak{}i\allowbreak{}l\allowbreak{}o\allowbreak{}n\allowbreak{}/\allowbreak{}(\allowbreak{}1\allowbreak{}-\allowbreak{}e\allowbreak{}p\allowbreak{}s\allowbreak{}i\allowbreak{}l\allowbreak{}o\allowbreak{}n\allowbreak{})\allowbreak{}.\allowbreak{}}
Both rank and physical time factors are explicit.

The same construction also supplies the spatial limit at fixed a,g with
xi\textless1/55. Every original Taylor coefficient of v has a connected
original face support, and a degree-n marked heat coefficient depends
only on its finite n-neighborhood. This is proved by the actual
source/heat recurrences: a derivative product joins supports only at an
original common edge; K inversion and its unperturbed resolvent preserve
link support. The scalar and disconnected terms remain until the
connected subtraction. On disjoint active components the original
compact-group operator and positive vacuum factor, giving the same zero
coefficients.

On two boxes containing the same N-neighborhood of the marks, all heat
coefficients through N coincide. Equation H13 consequently gives

\begingroup\small\begin{equation*}
\adjustbox{max width=0.92\linewidth}{$\displaystyle |\widehat C_L-\widehat C_{L'}|
\le2Ce^{-d\tau}\frac{(55|\xi|)^{2(\lfloor N/2\rfloor+1)}}{1-(55|\xi|)^2}.$}\tag{H34}
\end{equation*}\endgroup

The mean H2 and full heat H7 have analogous local Taylor tails. They
converge on cylinder Fourier polynomials, and for real coupling their
actual Markov contractions extend convergence to the uniform closure.
Positivity, unit mass, semigroup composition and symmetry pass by
cylinder approximation. Their unique invariant physical state follows by
integrating the uniform decay of H7 on centered cylinders. This is
uniqueness on the gauge-invariant observable algebra; the original
vertex Haar average supplies its gauge-invariant probability lift.

The resulting self-adjoint generator on L2(mu)\_phys has its actual
cylinder domain and A\_infinity
{F\allowbreak{}=\allowbreak{}k\allowbreak{}a\allowbreak{}p\allowbreak{}p\allowbreak{}a\allowbreak{}(\allowbreak{}K\allowbreak{}-\allowbreak{}D\allowbreak{}\_\allowbreak{}v\allowbreak{},\allowbreak{}i\allowbreak{}n\allowbreak{}f\allowbreak{}i\allowbreak{}n\allowbreak{}i\allowbreak{}t\allowbreak{}y\allowbreak{})\allowbreak{}F\allowbreak{}:\allowbreak{}}
pass the finite semigroup integral equation using the absolute local
derivative tails. The finite gap passes by L2 density, yielding
A\_infinity\textgreater=kappa d5(xi) on the centered space. The moment
matrices are bounded positive operators on the original l2 infinite face
coefficient space, because the row estimates and spatial tails give
strong convergence on finite coefficient vectors, then on all l2. Their
positive lower bounds remain.

Consequently R\_col and Phi extend as closed-range maps and A\_infinity
Phi=kappa R\_col. In particular A\_infinity P\_Phi is bounded even for
this infinite family. The same bounded-off-diagonal domain argument in
H25, and the complete H26 minimization, therefore hold for the full
infinite complementary physical space. The benchmark table H5 carries
over at fixed spacing and coupling.

\subsection{H7. Continuum scope and next original
quantity}\label{h7.-continuum-scope-and-next-original-quantity}

The standing simultaneous path keeps a\_n=a0*2\^{}(-n), g\_n²=1/c\_n,
{c\allowbreak{}\_\allowbreak{}n\allowbreak{}=\allowbreak{}g\allowbreak{}0\allowbreak{}\textasciicircum{}\allowbreak{}(\allowbreak{}-\allowbreak{}2\allowbreak{})\allowbreak{}+\allowbreak{}b\allowbreak{}e\allowbreak{}t\allowbreak{}a\allowbreak{}*\allowbreak{}n\allowbreak{}*\allowbreak{}l\allowbreak{}o\allowbreak{}g\allowbreak{}2\allowbreak{}}
and xi\_n=c\_n²/4. Its exact inclusion in the fifth-source domain is
{c\allowbreak{}\_\allowbreak{}n\allowbreak{}\textless{}\allowbreak{}=\allowbreak{}2\allowbreak{}s\allowbreak{}q\allowbreak{}r\allowbreak{}t\allowbreak{}(\allowbreak{}a\allowbreak{}l\allowbreak{}p\allowbreak{}h\allowbreak{}a\allowbreak{})\allowbreak{};\allowbreak{}}
its inclusion in this heat circle is c\_n\textless2/sqrt(55). For
beta\textgreater0 it eventually leaves both. No positive finite
continuum mass or nontrivial smooth four-dimensional field is concluded
from these fixed-spacing estimates.

This tranche completes the selected fifth-reference correction and its
return to the full physical complement. Merely re-centering the
elementary norm inequality reproduces its old endpoint, as the preserved
ROUTE\_ASSESSMENT calculation proves. A further extension has to use
actual additional coefficient or signed inverse information. The
concrete next candidate is the same-support linearized action J5 and its
preconditioned residual, using the now stored complete edge-channel
coefficients; any replacement of the inverse must retain its exact
residual operator. The missing sixth-source catalogue remains missing in
the delivered source chain and has not been declared reconstructed by
these estimates.

\subsection{H8. Evidence and
attribution}\label{h8.-evidence-and-attribution}

All old heat coefficients and absolute support sums are preserved
byte-for-byte. New algebra evaluates the fifth reference; new rational
calculations verify its root, heat circle, signed time integral and the
benchmark matrix inequalities. The written source-domain and semigroup
argument is supplied in F1-\/-F42 and H1-\/-H34; exact finite checks are
not represented as formal proof of the analytic passage.

G. Lumer and R. S. Phillips, Dissipative operators in a Banach space,
Pacific J. Math.11 (1961),679-\/-698, Theorem3.1, is used only after the
domain, dissipativity and range checks H5. P. Eymard\textquotesingle s
compact Fourier algebra and Schuette-\/-Zheng-\/-Hamer\textquotesingle s
exp-S/Casimir construction are credited as in the companion. This
contribution is on the original Yang--Mills objects. No arithmetic Gamma
measure, RH root packet, or Riemann activation asymptotic is imported.
\end{document}
